\newpage
\chapter{矩阵}
\setcounter{section}{0}
\section{从线性映射看矩阵与行列式}
前面我们并未纠结矩阵的引入问题,这里略引
线性映射来作解释.二者的深刻关系需要到后期才会被
解释.下文只考虑实矩阵.


我们考虑全体$n$阶方阵的集合为$
\bm{M}_n\left(\mathbb{R} 
\right)$,然后构造这样一个映射$det$:
\begin{align*}
   \det : \bm{M}_n\left(\mathbb{R} 
   \right)  & \longrightarrow 
\mathbb{R} \\
\bm{A} & \longmapsto 
\left|\bm{A}\right|
\end{align*}
即函数:
\[
\left|\bm{A}\right| =  \det\,\left(\bm{A}\right)
\]
这样的线性映射就是矩阵和行列式的内在关系.


说白了,我们的研究,就是为了回答下面这两个问题:
\begin{enumerate}[label =\textup{(\arabic*)}]
\item $\det\,\left(\bm{A}\right)$能够反映出$\bm{A}$的哪些性质？
\item 映射$\det$本身的性质如何？
\end{enumerate}

进一步地,我们将要探讨代数与几何二者之间的关系.
\newpage
\section{矩阵的基本形式}
\subsection{矩阵的表示}
\begin{formal}
\begin{definition}[矩阵的形式]\label{def:矩阵的形式}
矩阵是一张\emph{数表},通常用大写英文字母表示,如:
\[  
    \bm{A} = \begin{pmatrix}
        a_{11}  & \dots & a_{1n} \\
        \vdots & \ddots  & \vdots \\
        a_{m1}  & \dots & a_{mn}
    \end{pmatrix}_{m \times n}
\]

这是一个$\displaystyle m \times n$的矩阵,即它有$\displaystyle 
m$行$\displaystyle n$列,其中,$\displaystyle a_{ij}$表示
第$\displaystyle i$行第$\displaystyle j$列的元素.$\forall 1\leqslant i\leqslant m,1\leqslant j\leqslant n,a_{ij}\in \mathbb{R}\left(\mathbb{C}\right)$,则称为一个$m\times n$阶的实（复）矩阵,记作$\bm{A}\in M_{m\times n}\left(\mathbb{R}\right)\left(M_{m\times n}\left(\mathbb{C}\right)\right).$

这个矩阵可记作$\displaystyle a_{ij}$或$\displaystyle \left(
    a_{ij}
\right)_{m \times n}$或$\displaystyle \bm{A}_{mn}$或$
\bm{A}_{m \times n}$.一些情况下也会记作
$\bm{A}^{m\times n}$等形式.
\end{definition}
\end{formal}
\begin{formal}
\begin{definition}[零矩阵]\label{def:零矩阵}
特别地,所有元素为$\displaystyle 0$的矩阵称为零矩阵,记作
$\displaystyle \bm{O}$
\[
    \bm{O} = \begin{pmatrix}
        0 &\cdots &0 \\
        \vdots& \ddots &\vdots\\
        0&\dots&0
    \end{pmatrix}
\]
\end{definition}
\end{formal}

任何一个矩阵中的$\displaystyle 0$都可以用空白代替.

\begin{formal}
\begin{definition}[负矩阵]\label{def:负矩阵}
    所谓$\displaystyle \bm{A}$的负矩阵就是把其中每个元素变为相反数:
\[  \displaystyle
    \bm{-A} = \begin{pmatrix}
        -a_{11}  & \dots & -a_{1n} \\
        \vdots & \ddots  & \vdots \\
        -a_{m1}  & \dots & -a_{mn}
    \end{pmatrix}_{m \times n}
\]
\end{definition}
\end{formal}
\begin{red}
\begin{remark}
    负矩阵也可以用\cref{def:数乘}来定义.
\end{remark}
\end{red}
\begin{red}
    \begin{remark}
        矩阵不是一个数.
    \end{remark}
\end{red}
\subsection{几种基础矩阵}
\begin{formal}
\begin{definition}[方阵]\label{def:方阵}
一个行数和列数相等的矩阵称为方阵,比如
\[ 
    \bm{A}_{n} = \begin{pmatrix}
        a_{11}  & \dots & a_{1n} \\
        \vdots & \ddots  & \vdots \\
        a_{n1}  & \dots & a_{nn}
    \end{pmatrix}_{n \times n}
\]
其中,$a_{11}$到$a_{nn}$这
一条连线称为主对角线.
\end{definition}
\end{formal}
\begin{formal}
\begin{definition}[单位阵]\label{def:单位阵}
如果一个方阵主对角线上全是$1$,其他元素都是$
 0$,那么它是一个单位矩阵
\[
    \bm{I}_{n}  = \bm{E}_n = \begin{pmatrix}
        1 & & \\
        & \ddots &\\
        &&1
    \end{pmatrix}
    _{n}
\]
\end{definition}
\end{formal}
\begin{formal}
\begin{definition}[常量矩阵]\label{def:常量矩阵}
如果一个方阵主对角线上全是一个相同的常数$k$,
其他元素都是零,称为数量矩阵或常量矩阵
\[
    \displaystyle
    \bm{K}_n = \begin{pmatrix}
        k & & \\
        & \ddots &\\
        &&k
    \end{pmatrix}
    _n
\]
\end{definition}
\end{formal}
\begin{formal}
\begin{definition}[对角矩阵]\label{def:对角矩阵}
如果一个方阵除了主对角线上其他元素都是零（主对角线上
不一定没有零）,称为对角矩阵
\[
    \bm{\varLambda} _n = \begin{pmatrix}
        a_{11} & & \\
        & \ddots &\\
        &&a_{nn}
    \end{pmatrix}= \mathrm{diag}\,\left\{
        a_{11},\cdots,a_{nn}
    \right\}
\]
\end{definition}
\end{formal}
\begin{formal}
\begin{definition}[三角矩阵]

如果一个方阵主对角线以上（下）全是$\displaystyle 0$
,称为下（上）三角阵
\[
    \bm{A}_n = \begin{pmatrix}
        a_{11} & a_{12} &\dots & a_{1n} \\
        & a_{22} & \dots & a_{2n}\\
        & & \ddots & \vdots \\
        & & & a_{nn}
    \end{pmatrix} 
\]

\[
    \displaystyle
    \bm{B}_n = \begin{pmatrix}
        b_{11} & & & \\
        b_{21} & b_{22} & & \\
        \vdots & \vdots & \ddots & \\
        b_{n1} & b_{n2} & \dots & b_{nn}
    \end{pmatrix}
\]
\end{definition}
\end{formal}
\begin{formal}
\begin{definition}[旋转矩阵]\label{def:旋转矩阵}
这是一种应用广泛而重要的矩阵,这里给出二阶形式,注意,逆时针为正方向.
\[
\bm{A} = \begin{pmatrix}
    \cos \theta  & -\sin\theta \\
    \sin\theta & \cos\theta 
\end{pmatrix}
= \mathrm{exp}\left(
    \theta \begin{bmatrix}
        0 & -1 \\
        1 & 0
    \end{bmatrix}
\right)
\]\end{definition}\end{formal}
\begin{red}
    \begin{remark}
        旋转矩阵对于计算机方向有重大意义.
    \end{remark}
\end{red}
\begin{formal}
\begin{definition}[阶梯形矩阵]\label{def:阶梯形矩阵}
设$\bm{A}_{mn}$
为非零矩阵,若所有非零行（即至
少有一个非零元素的行）全在零行的上
面,$\bm{A}_{mn}$中各非零行中第一个
（最后一个）非
零元素（称为阶
梯点）前（后）面零元素的个数随行数增
大而严格增多（减少）,则称为上（下）梯
形矩阵,简称为上（下）梯形阵,上下梯形阵
统称为梯形阵.
\end{definition}
\end{formal}

\begin{red}
\begin{remark}
    一个矩阵有可能同时是上、下梯形阵,也可能都不是,也可能
只是其中之一！
\end{remark}
\end{red}
\subsection{矩阵的迹}
\begin{formal}
\begin{definition}[迹]\label{def:迹}
方阵
\[
\begin{pmatrix}
       a_{11} &a_{12}&\cdots&a_{1n}\\
         a_{21} &a_{22}&\cdots&a_{2n}\\
            \vdots&\vdots&&\vdots\\
            a_{n1}&a_{n2}&\cdots&a_{nn} 
\end{pmatrix}
\]
的主对角线上元素之和称为这个方阵的迹,记作$ \mathrm{tr}\,\bm{A}=\mathrm{Tr}\,\left(\bm{A}\right)=
  a_{11} + a_{22} + \cdots + a_{nn}$.
\end{definition}
\end{formal}
\begin{formal}
    \begin{theorem}[迹的交换性]\label{thm:迹的交换性}
        设$\bm{A}\in M_{m\times n}\left(\mathbb{R}\right),\bm{B}\in M_{n\times m}\left(\mathbb{R}\right)$,则\[
        \mathrm{Tr}\,\left(\bm{AB}\right)=\mathrm{Tr}\,\left(\bm{BA}\right)
        \]
    \end{theorem}
\end{formal}
\newpage 
\section{矩阵的运算}
\subsection{矩阵的线性运算}
\begin{formal}
\begin{definition}[同型]\label{def:同型}
    矩阵的同型指的是他们的行数和列数分别相等,如$
\bm{A}_{m\times n}$和$\bm{B}_{m\times n}$.
\end{definition}
\end{formal}
\begin{formal}
\begin{definition}[加法和减法]\label{def:加法和减法}
矩阵的加减法只能对同型矩阵做,对应位置相加减即可:
\[
    \left(a_{ij}\right)_{m \times n} \pm \left(
        b_{ij}
    \right)_{m \times n} = \left(
        a_{ij} \pm b_{ij}
    \right)_{m \times n}
\]
\[  \displaystyle
    \begin{pmatrix}
        a_{11}  & \dots & a_{1n} \\
        \vdots & \ddots  & \vdots \\
        a_{m1}  & \dots & a_{mn}
    \end{pmatrix}_{m \times n}
    \pm
    \quad
    \begin{pmatrix}
        b_{11}  & \dots & b_{1n} \\
        \vdots & \ddots  & \vdots \\
        b_{m1}  & \dots & b_{mn}
    \end{pmatrix}_{m \times n}
   = \quad
    \begin{pmatrix}
        a_{11} \pm b_{11} & \dots & a_{1n} \pm b_{1n} \\
        \vdots & \ddots  & \vdots \\
        a_{m1} \pm b_{m1}  & \dots & a_{mn}  \pm b_{mn}
    \end{pmatrix}_{m \times n}
\]
\end{definition}
\end{formal}
\begin{formal}
\begin{proposition}[加减法性质]\label{prop:加减法性质}
    矩阵的加减法满足:
\setcounter{equation}{0}
\begin{enumerate}[label=\textup{(\arabic*)}]
    \item 交换律：$\bm{A} + \bm{B} = \bm{B} + \bm{A} $
    \item 结合律：$\left( \bm{A} + \bm{B} \right) + \bm{C} = \bm{A}
    + \left( \bm{B} + \bm{C}\right)$
    \item 零元存在：$\bm{A} + \bm{O} = \bm{A}$
    \item 存在负矩阵：$
    \bm{A}+\left(-\bm{A}\right) = \bm{O}$
    \item 移项：$
    \bm{A} + \bm{B} = \bm{C} \Longleftrightarrow \bm{A}
    = \bm{C} - \bm{B} $
\end{enumerate}
\end{proposition}
\end{formal}
\subsection{矩阵的数乘}
\begin{formal}
\begin{definition}[数乘]\label{def:数乘}
一个常数$k$乘一个矩阵$ 
\bm{A}_{mn}$称为数乘运算,所有元素都乘这个常数即可:
\[
    k\bm{A}_{mn} = k\left(a_{ij}\right)_{m \times n}
     =\left(k a_{ij}\right)_{m \times n}
\]
\[  
    k \begin{pmatrix}
        a_{11}  & \dots & a_{1n} \\
        \vdots & \ddots  & \vdots \\
        a_{m1}  & \dots & a_{mn}
    \end{pmatrix}_{m \times n}
    = \quad
     \begin{pmatrix}
        ka_{11}  & \dots & ka_{1n} \\
        \vdots & \ddots  & \vdots \\
        ka_{m1}  & \dots & ka_{mn}
    \end{pmatrix}_{m \times n}
\]
\end{definition}
\end{formal}
\begin{formal}
\begin{proposition}[数乘性质]\label{prop:数乘性质}
矩阵的数乘满足:
\begin{enumerate}[label=\textup{(\arabic*)}]
   \item 数的分配律$k \left(\bm{A} + \bm{B}\right) = k\bm{A} + k\bm{B}$
   \item 矩阵的分配律$\left(k + l\right)\bm{A} = k\bm{A} + l \bm{B}$
   \item 结合律$k\left(l\bm{A}\right) = 
   \left(kl\right) \bm{A} = 
   l \left(k\bm{A}\right)$
    \item 单位元存在$1\cdot\bm{A} = \bm{A}$
    \item 零元存在$0\cdot\bm{A} = \bm{O}$
\end{enumerate}
\end{proposition}
\end{formal}
\subsection{矩阵的乘法}
矩阵的乘法首先要满足这样的条件:前一个矩阵的列数与后一个矩阵的
行数相等.而它们相乘的结果是一个行数等于前一个矩阵的行数,列数
等于后一个矩阵的列数的矩阵,简单说就是:“中间相等,取两头”.
\[
    \bm{A}_{m \times n} \bm{B}_{n \times k} = \bm{C}_{m \times k}
\]
\begin{formal}
\begin{definition}[矩阵乘法]\label{def:矩阵乘法}
    对于矩阵乘法$\bm{A}_{m\times n}\bm{B}_{n\times k}=\bm{C}_{m\times k}$结果矩阵中的元素$c_{ij}$的算法是用矩阵$
\bm{A}$中的第$i$行$\left(a_{i1},\dots ,a_{in
}\right)$各元素分别与矩阵$\bm{B}$中的第$
j$列$\left(b_{1j},\dots,
b_{nj}\right)$各元素分别相乘,再把结果相加得
到$c_{ij}$即这两个向量做内积；
\[
c_{ij}=
    \begin{pmatrix}
        a_{i1} & \dots &a_{in}
    \end{pmatrix}
    \cdot 
    \begin{pmatrix}
        b_{1j}\\
        \vdots\\
        b_{nj}
    \end{pmatrix}\Longleftrightarrow
    c_{ij} = \sum _{k=1} ^{n} a_{ik}b_{kj}
\]
\end{definition}
\end{formal}
\begin{red}
\begin{remark}
    矩阵乘法最大的特点就是不能随意交换即
\[
 \bm{A}\bm{B}=\bm{B}\bm{A}
\]
多数情况下不成立.
\end{remark}
\end{red}
 \begin{formal}
 \begin{proposition}[矩阵乘法的性质]\label{prop:矩阵乘法的性质}
    设$\bm{A}_{m \times n} =
 \left(a_{ij}\right)_{m \times n}$,
 $\bm{B}_{n \times p} = \left(
    b_{ij}
 \right)_{n \times p}$,$\bm{C}
 _{p \times q} = \left(
    c_{ij}
 \right)_{p \times q}$,则
\begin{enumerate}[label=\textup{(\arabic*)}]
    \item 结合律：$\left(\bm{A}\bm{B}\right)\bm{C} = \bm{A}\left(\bm{B}\bm{C}\right) $
    \item 分配律：$\bm{A}\left(\bm{B} + \bm{C}\right) = \bm{A}\bm{B} + \bm{A}\bm{C} $
    \item 分配律：$\left(\bm{B}+\bm{C}\right)\bm{A}=\bm{B}\bm{A}+\bm{C}\bm{A} $
    \item 与数乘相容：$k\left(\bm{A}\bm{B}\right) = 
    \left(k\bm{A}\right)\bm{B}
     = \bm{A}\left(k\bm{B}\right)$
    \item 单位元：$
    \bm{E}_m \bm{A}_{m \times n} =
     \bm{A}_{m \times n} =
     \bm{A}_{m \times n}\bm{E}_{n} $
    \item 注意:$\left(\bm{A}+\bm{B}\right)^2=
    \bm{A}^2+\bm{A}\bm{B}+\bm{B}
    \bm{A}+\bm{B}^2.$
    \item 对于方阵$\bm{A}\text{、}
    \bm{B}, \det\left(
        \bm{AB}
    \right) =  \det\left(\bm{A}\right)
    \cdot  \det\left(\bm{B}\right)$
\end{enumerate}
 \end{proposition}
 \begin{Proof}

 \begin{enumerate}[label=\textup{(\arabic*)}]
    \item 因为$\bm{AB}$的$\left(i,j\right)$
 元为
 \[
 \sum_{k = 1}^{n}a_{ik}b_{kj}
 \]
 $\left(\bm{AB}\right)\bm{C}$的$
 \left(i,j\right)$元为
 \[
 \sum_{r = 1}^{p}\left(\sum_{k = 1}
 ^{n}a_{ik}b_{kr}\right)c_{kj}
 =
 \sum_{r = 1}^{p}\sum_{k = 1}
 ^{n}a_{ik}b_{kr}c_{kj}
 \]
 $\bm{A}\left(\bm{BC}\right)$的
 $\left(i,j\right)$元为
 \[
 \sum_{k = 1}^{n}\sum_{r = 1}
 ^{p}a_{ik}b_{kr}c_{kj}\qedhere
 \]
 \end{enumerate}
 \end{Proof}
 \end{formal}
\begin{red}
\begin{remark}
    简单说一下左乘和右乘,矩阵$\bm{A}$左乘矩阵$\bm{B}$指的是矩阵$\bm{A}$从左乘矩阵$\bm{B}$即得$\bm{AB}$.
\end{remark}
\end{red}
\begin{red}
\begin{remark}
    一般而言,因为整性不成立,矩阵的乘法消去律亦不成立,
    从线性映射的角度来看,两个非零映射符合得到零映射也并非不可能.
\end{remark}
\end{red}
\subsection{方阵的非负整数次幂}
\begin{formal}
\begin{definition}[方阵的幂]\label{def:方阵的幂}
对于$n$阶方阵$\bm{A}_{n}$,定义它的
非负整数$k$次幂为:
\begin{equation*}
    \left(\bm{A}_n\right)^k = 
    \begin{cases}
        \bm{A}\bm{A} \cdots \bm{A} ,& k \in \mathbb{N} ^*;\\
        \bm{E}_n,&k = 0 .
    \end{cases}
\end{equation*}
\end{definition}
\end{formal}
\begin{formal}
\begin{proposition}[幂的性质]\label{prop:幂的性质}
    \begin{enumerate}[label=\textup{(\arabic*)}]
    \item$\bm{A}^{k + l} = \bm{A}^k\bm{A}^l$
    \item$\bm{A}^{kl} = \left(\bm{A}^{k}\right)^{l}$
    \item 由于矩阵乘法的不可交换,一般有:$\left(
        \bm{A}\bm{B}
    \right)^k\neq \bm{A}^k \bm{B}^k$
\end{enumerate}
\end{proposition}
\end{formal}
\begin{formal}
\begin{definition}[方阵的多项式]\label{def:方阵的多项式}
    类似实数的多项式:
\[
    f\left(x\right) = a_nx^n + \cdots a_1x+a_0.
\]
定义方阵$\bm{A}_n$的多项式:
\[
f\left(\bm{A}_n\right) = a_n\left(\bm{A}_n\right)^n + \cdots a_1\bm{A}_n+a_0\bm{E}_n.
\]
\end{definition}
\end{formal}
\begin{red}
    \begin{remark}
        多项式是后续学习的一大重心,它是重要的结构.
    \end{remark}
\end{red}
\newpage
\section{矩阵的转置}
\subsection{转置}
\begin{formal}
\begin{definition}[转置]
    矩阵的转置是将它的行和列交换,对\[
\displaystyle
     \bm{A} = \begin{pmatrix}
        a_{11} & \cdots &a_{1n}\\
        \vdots &\ddots &\vdots\\
        a_{m1} &\cdots &a_{mn}
    \end{pmatrix}_{m \times n}
\]
其转置
\[
    \bm{A}' =
    \bm{A}^{\prime} = \begin{pmatrix}
        a_{11} & \cdots &a_{m1}\\
        \vdots &\ddots &\vdots\\
        a_{1n} &\cdots &a_{mn}
    \end{pmatrix}_{n\times m}
\]
\end{definition}
\end{formal}
\begin{formal}
\begin{proposition}[转置的性质]\label{prop:转置的性质}矩阵的转置满足：

    \begin{enumerate}[label=\textup{(\arabic*)}]
    \item$\left(\bm{A}^{\prime}\right)^{\prime} = \bm{A}$
    \item$\left(\bm{A}+\bm{B}\right)^{\prime} = \bm{A}^{\prime} + \bm{B}^{\prime}$
    \item $\left(k\bm{A}\right)^{\prime} = k\bm{A}^{\prime}$
    \item 但是,最重要的是:
    $\left(\bm{A}\bm{B}\right)^{\prime} = \bm{B}^{\prime}\bm{A}^{\prime}$
    \item 特别地:$
    \bm{\varLambda} '= 
    \bm{\varLambda} $ 
\end{enumerate}
\end{proposition}
\end{formal}
\subsection{对称阵与反对称阵}
\begin{formal}
\begin{definition}[对称阵与反对称阵]\label{def:对称阵与反对称阵}
如果一个方阵满足$\displaystyle \bm{A}^{\prime} = \bm{A}$,那么它是
一个对称阵；如果满
足$\displaystyle
 \bm{A}^{\prime} = -\bm{A}$,那么它是
 一个反对称阵
\begin{align*}
    & \text{对称阵：}\bm{A}^{\prime} =
     \bm{A} \Longleftrightarrow   
    a_{ij}=a_{ji};\\
    & \text{反对称阵：} 
    \bm{A}^{\prime} = \bm{A}\Longleftrightarrow 
     a_{ij}
    =-a_{ji}\text{且}a_{ii}=0.
\end{align*}
\end{definition}
\end{formal}
\begin{red}\begin{remark}
这些方阵一定是对称阵:\[
    \bm{A}\bm{A}^{\prime},\bm{A}^{\prime}\bm{A},\bm{A}+\bm{A}^{\prime}.
\]
\end{remark}
\end{red}
\begin{formal}
\begin{theorem}[矩阵拆分]\label{thm:矩阵拆分}
    一般地,任何一个方阵都可以分解为一
    个对称阵和一个反对称阵即
\[
    \bm{A} = \dfrac{\bm{A}+\bm{A}^{\prime}}{2}  + \dfrac{\bm{A}
    -\bm{A}^{\prime}}{2}.
\]
\end{theorem}
\end{formal}
\begin{formal}
    \begin{proposition}[对称与反对称刻画]\label{prop:对称与反对称刻画}
        设$\bm{A}$、$\bm{B}$为$n$阶方阵,则
        \begin{enumerate}[label=\textup{(\arabic*)}]
            \item 若$\bm{A},\bm{B}$为对称阵,
            则$\bm{AB}$为对称阵等价于$\bm{AB}=\bm{BA}$
            ；$\bm{AB}$为反对称阵等价于
            $\bm{AB}=-\bm{BA}.$
            \item 若$\bm{A}$为对称阵,
            $\bm{B}$为反对称阵,
            则$\bm{AB}$为反对称阵等价于$\bm{AB}=\bm{BA}$
            ；$\bm{AB}$为对称阵等价于
            $\bm{AB}=-\bm{BA}.$
        \end{enumerate}
    \end{proposition}
    \begin{Proof}
    给出若$\bm{A},\bm{B}$为对称阵,
            则$\bm{AB}$为对称阵等价于$\bm{AB}=\bm{BA}$的证明,
            其余同理.

        一方面,因为$\bm{AB}$为对称阵,故
        \[
        \bm{AB}=\left(\bm{AB}\right)^{\prime}=\bm{B}^{\prime}\bm{A}^{\prime}=\bm{BA}
        \]
        另一方面,因为$\bm{AB}=\bm{BA}$,故
        \[
        \left(\bm{AB}\right)^{\prime}=\bm{B}^{\prime}\bm{A}^{\prime}=\bm{BA}=\bm{AB}
        \qedhere \]
    \end{Proof}
\end{formal}
对于复矩阵,也有一种类似对称的定义.
\begin{formal}
\begin{definition}[Hermite矩阵]\label{def:Hermite矩阵}
    设$n$阶复矩阵满足$\overline{\bm{A}}^{\prime}=\bm{A}$,
    则称为一个\textup{Hermite}（厄米特）矩阵；
    若$\overline{\bm{A}}^{\prime}=-\bm{A}$,称为是斜
    \textup{Hermite}矩阵.
\end{definition}
\end{formal}
\begin{formal}
\begin{theorem}[复矩阵拆分]\label{thm:复矩阵拆分}
    任一$n$阶复矩阵均可表示为一个
    \textup{Hermite}矩阵与一个
    斜\textup{Hermite}矩阵之和即
    \[
    \bm{A}=\frac{\bm{A}+\overline{\bm{A}}^{\prime}}{2}+
    \frac{\bm{A}-\overline{\bm{A}}^{\prime}}{2}
    \]
\end{theorem}
\end{formal}
\newpage
\section{复矩阵的共轭}
\begin{formal}
\begin{definition}[共轭]\label{def:共轭}
    对于复矩阵,其矩阵的共轭等于元共轭的矩阵
\[
\overline{\bm{A}} =
\left(\overline{a_{ij}} \right)
_{m \times n}
\]
\end{definition}
\end{formal}
\begin{formal}
\begin{proposition}[共轭的性质]\label{prop:共轭的性质}
    共轭满足的性质有：
    \begin{enumerate}[label=\textup{(\arabic*)}]
    \item $\displaystyle \overline{\left(\overline{\bm{A}}\right)}=\bm{A}$
    \item $\overline{\bm{A} + \bm{B}} = 
    \overline{\bm{A}} + 
    \overline{\bm{B}}  $
    \item 若$c \in \mathbb{C}$
        ,则
    $\overline{c\bm{A}} = \overline{c}
    \overline{\bm{A}}$
    \item $\overline{\bm{AB}} =
    \overline{\bm{A}}
    \overline{\bm{B}}$
    \item $\overline{
        \left(\bm{A}^{\prime}\right)}
    = \left(\overline{\bm{A}} 
    \right)^{\prime}$ 
    \item $\left| \overline{\bm{A}}\right|  =
    \overline{\left|\bm{A}\right|}$（利用组合定义非常容易证明）
\end{enumerate}
\end{proposition}
\end{formal}
\newpage
\section{方阵的逆阵}
\subsection{定义}
\begin{formal}
\begin{definition}[逆阵]\label{def:逆阵}
    如果对于一个$n$阶方阵$\bm{A}$,存在一个
    $n$阶方阵$\bm{B}$使得
    \[
    \bm{AB} = \bm{BA} = \bm{I}_n
    \]
    那么$\bm{B}$称为$\bm{A}$的逆矩阵,记作
    \[
    \bm{B} = \bm{A}^{-1}
    \]
\end{definition}
\end{formal}
\begin{red}
\begin{remark}
    只有方阵才具有逆阵.
\end{remark}
\end{red}
\begin{red}
\begin{remark}
    并不是所有非零方阵都具有逆阵.
\end{remark}
\end{red}
\begin{formal}
    \begin{definition}[奇异阵]\label{def:奇异阵}
称有逆阵的方阵为可逆阵或非奇异阵或非异阵,
而称无逆阵的方阵为奇异阵.
\end{definition}
\end{formal}
\subsection{性质}
\begin{formal}
\begin{proposition}[逆阵的性质]\label{prop:逆阵的性质}
    逆阵的性质有（以下均默认可逆）：
\begin{enumerate}[label=\textup{(\arabic*)}]
    \item 逆阵存在则必定唯一
    \item $\left(
        \bm{A}^{-1}
        \right)^{-1} = \bm{A}$
    \item $\bm{A}$、$
    \bm{B}$均可逆则$
    \bm{AB}$可逆且$
    \left(\bm{AB}\right)^{-1} = 
    \bm{B}^{-1}\bm{A}^{-1}$
    \item 若$0\neq c \in \mathbb{R}
    $,则
     $c\bm{A}$可逆且
     $\left(c\bm{A}\right)^{-1}=
     c^{-1}\bm{A}^{-1}$
    \item $\bm{A}$
    可逆则$\bm{A}^{\prime}$
    可逆且$
    \left(\bm{A}^{\prime}\right)^{-1} 
    = \left(
    \bm{A}^{-1}
    \right)^{\prime}$
    \item 一般$
    \bm{AB}\neq\bm{BA}$
        时
    $\left(
        \bm{AB}\right)^{-1}
    \neq\bm{A}^{-1}
    \bm{B}^{-1}$
    \item 整性对可逆阵成立
    \item 对可逆阵,乘法消去律成立
\end{enumerate}
\end{proposition}
\end{formal}
\begin{green}
\begin{corollary}[逆阵的积]\label{cor:逆阵的积}
    设可逆阵$\bm{A}_1,\bm{A}_2,\cdots,\bm{A}_m\in M_n\left(\mathbb{R}\right)$,则$\bm{A}_1\bm{A}_2\cdots\bm{A}_m$也可逆且\[
    \left(
        \bm{A}_1\bm{A}_2\cdots\bm{A}_m  
    \right)^{-1}= \bm{A}_m^{-1}\bm{A}_{m-1}^{-1}\cdots\bm{A}_1^{-1}
    \]
\end{corollary}
\end{green}
\subsection{伴随阵}
\begin{formal}
\begin{definition}[伴随阵]\label{def:伴随阵}
    $\bm{A}$为$n$阶方阵,$A_{ij}$为$ \det\left(
        \bm{A}
    \right)$的$\left(i,j\right)$元的代数余子式
    ,则$\bm{A}$的伴随阵为
    \[
    \bm{A}^* = \begin{pmatrix}
        A_{11} & \cdots & A_{n1} \\
        \vdots & \ddots & \vdots \\
        A_{1n} & \cdots & A_{nn}
    \end{pmatrix}
    \]
\end{definition}
\end{formal}
\begin{red}
\begin{remark}
  关于下标,应当以$ \det\left(\bm{A}\right)$
剩下的元的位置看而不是划去的元的位置.或者也可以说多了一步转置.  
\end{remark}
\end{red}
\begin{red}
\begin{remark}
    矩阵的伴随矩阵一定存在.
\end{remark}
\end{red}
\begin{green}
\begin{lemma}[与伴随阵的积]\label{lem:与伴随阵的积}
    设$n$阶方阵$\bm{A}$的伴随阵$\bm{A}^*$,则由
    异乘变零定理可得
    \[
    \bm{A}\bm{A}^* =\bm{A}^*\bm{A}
      = \left|\bm{A}\right|\bm{I}_n
     =  \det\left(\bm{A}\right)\bm{I}_n
    \]
\end{lemma}
\begin{Proof}
考虑到\cref{thm:异乘变零定理}异乘变零定理,有
\begin{align*}
    \begin{vmatrix}
        a_{11}&a_{12}&\cdots&a_{1n}\\
        a_{21}&a_{22}&\cdots&a_{2n}\\
        \vdots&\vdots&&\vdots\\
        a_{n1}&a_{n2}&\cdots&a_{nn}
    \end{vmatrix}\begin{vmatrix}
        A_{11}&A_{21}&\cdots&A_{n1}\\
        A_{12}&A_{22}&\cdots&A_{n2}\\
        \vdots&\vdots&&\vdots\\
        A_{1n}&A_{2n}&\cdots&A_{nn}
    \end{vmatrix}=
    \begin{vmatrix}
        \left|\bm{A}\right|&0&\cdots&0\\
        0&\left|\bm{A}\right|&\cdots&0\\
        \vdots&\vdots&&\vdots\\
        0&0&\cdots&\left|\bm{A}\right|
    \end{vmatrix}\qedhere
\end{align*}
\end{Proof}
\end{green}
\begin{formal}
\begin{theorem}[逆阵存在性]\label{thm:逆阵存在性}
    对于矩阵$\bm{A}$
    \[\det\,\left(\bm{A}\right)
    \begin{cases*}
    =0\Longrightarrow \text{不可逆/奇异}\\
    \neq 0\Longrightarrow \text{可逆/非奇异}
    \end{cases*}
    \]
\end{theorem}
\end{formal}
\begin{formal}
\begin{theorem}[逆阵与伴随]\label{thm:逆阵与伴随}
    若$\bm{A}$可逆,即$ \det\left(
        \bm{A}
    \right) \neq 0$,那么
    \[
    \bm{A}^{-1} = \frac{1}{\left|\bm{A}
    \right|} \bm{A}^*
    \]
\end{theorem}
\begin{Proof}
$\left|\bm{A}\right| \neq 0$时
\[
\bm{A}\left(\frac{1}{\left|\bm{A}\right|}\bm{A}^*\right)=
\frac{\bm{AA}^*}{\left|\bm{A}\right|}=
\bm{I}_n \qedhere
\]
\end{Proof}
\end{formal}
\begin{red}
\begin{remark}
    有了这样一套矩阵的逆阵理论,对于\cref{thm:Cramer法则}中的线性方程组
    \[
    \begin{cases*}
    a_{11}x_1+a_{12}x_2+\cdots+a_{1n}x_n=b_1\\
    a_{21}x_1+a_{22}x_2+\cdots+a_{2n}x_n=b_2\\
    \qquad\qquad\cdots\cdots\cdots\cdots\\
    a_{n1}x_1+a_{n2}x_2+\cdots+a_{nn}x_n=b_n
    \end{cases*}\Longleftrightarrow \bm{Ax}=\bm{\beta}
    \]在$\left|\bm{A}\right|\neq 0$时就有解为
    \[
    \bm{x}=\bm{A}^{-1}\bm{\beta}
    \]
\end{remark}
\end{red}
\begin{formal}
\begin{theorem}[方阵积的行列式]\label{thm:方阵积的行列式}
    $\bm{A}$、$\bm{B}$为$n$阶方阵则
    \[
    \left|\bm{AB}\right|=\left|\bm{A}\right|\left|\bm{B}\right|
    \]
\end{theorem}
\begin{Proof}
设$\bm{A}=\left(a_{ij}\right)_{n \times n}$,$\bm{B} = 
\left(b_{ij}\right)_{n \times n}$,构造
\[
\bm{C}_{2n \times 2n}=\begin{bmatrix}
    \bm{A}&\bm{O}\\
    -\bm{I}_n&\bm{B}
\end{bmatrix}
\]
则
\begin{align*}
 \left|\bm{C}\right|   &=
 \begin{vmatrix}
    a_{11}&a_{12}&\cdots&a_{1n}&0&0&\cdots&0\\
    a_{21}&a_{22}&\cdots&a_{2n}&0&0&\cdots&0\\
    \vdots&\vdots&\ddots&\vdots&\vdots&\vdots&\ddots&\vdots\\
    a_{n1}&a_{n2}&\cdots&a_{nn}&0&0&\cdots&0\\
    -1&0&\cdots&0&b_{11}&b_{12}&\cdots&b_{1n}\\
    0&-1&\cdots&0&b_{21}&b_{22}&\cdots&b_{2n}\\
    \vdots&\vdots&\ddots&\vdots&\vdots&\vdots&\ddots&\vdots\\
    0&0&\cdots&-1&b_{n1}&b_{n2}&\cdots&b_{nn}
 \end{vmatrix}\\
 &=
 \begin{vmatrix}
    0&0&\cdots&0&\sum\limits_{r}a_{1r}b_{r1}&\sum\limits_{r}a_{1r}b_{r2}&\cdots&\sum\limits_{r}a_{1r}b_{rn}\\
    0&0&\cdots&0&\sum\limits_{r}a_{2r}b_{r1}&\sum\limits_{r}a_{2r}b_{r2}&\cdots&\sum\limits_{r}a_{2r}b_{rn}\\
    \vdots&\vdots&\ddots&\vdots&\vdots&\vdots&\ddots&\vdots\\
    0&0&\cdots&0&\sum\limits_{r}a_{nr}b_{r1}&\sum\limits_{r}a_{nr}b_{r2}&\cdots&\sum\limits_{r}a_{nr}b_{rn}\\
    -1&0&\cdots&0&b_{11}&b_{12}&\cdots&b_{1n}\\
    0&-1&\cdots&0&b_{21}&b_{22}&\cdots&b_{2n}\\
    \vdots&\vdots&\ddots&\vdots&\vdots&\vdots&\ddots&\vdots\\
    0&0&\cdots&-1&b_{n1}&b_{n2}&\cdots&b_{nn}
 \end{vmatrix}\\
 &=
 \begin{vmatrix}
    \bm{O}&\bm{AB}\\
    -\bm{I}_n&\bm{B}
 \end{vmatrix}
\end{align*}
接下来用Laplace定理以前$n$行展开,化简前
\[
\left|\bm{C}\right|=
\begin{vmatrix}
    \bm{A}&\bm{O}\\
    -\bm{I}_n&\bm{B}
\end{vmatrix}=\left| \bm{A} \right|
 \left| \bm{B} \right|
\]
化简后
\begin{align*}
\left| 
    \bm{C}
    \right|
    &=
    \left( -1 \right) ^{1+2+\cdots+2n}
        \left|\bm{AB}\right|
        \left|-\bm{I}_n\right|\\
&=\left(-1\right)^{n^2+n}\left|\bm{AB}\right|\\
&=\left|\bm{AB}\right|
\qedhere
\end{align*}
\end{Proof}
\end{formal}
\subsection{方阵可逆的充要条件}
\begin{formal}
    \begin{theorem}[矩阵可逆的充要条件]
对于$\bm{A}$,
    \begin{equation*}
         \det\left(\bm{A}\right)
        \begin{cases*}
            = 0 \Longrightarrow\text{不可逆}\\
            \neq 0 \Longrightarrow \text{可逆}
        \end{cases*}
    \end{equation*}
\end{theorem}
\end{formal}
\begin{green}
\begin{corollary}[逆阵的若干推论]\label{cor:逆阵的若干推论}对于逆阵,有：

    \begin{enumerate}[label=\textup{(\arabic*)}]
        \item $\bm{A}_1,\cdots,\bm{A}_m$
            均为
        $n$阶方阵,且存在$\bm{A}_i$是
        奇异阵,那么$\bm{A}_1
        \cdots \bm{A}_m$是奇异阵
        \item $ \bm{A}$
        可逆$\Longleftrightarrow 
        \left|\bm{A}\right|^{-1} = \left|
            \bm{A}^{-1}
        \right|$
        \item 对于$n$阶方阵$
        \bm{A}$、$\bm{B}$,若$
        \bm{AB}=\bm{I}_n$或$
        \bm{BA}=\bm{I}_n\Longleftrightarrow 
        \bm{B} = \bm{A}^{-1}$
    \end{enumerate}
\end{corollary}
\begin{Proof}
$(1)$计算行列式即证.

$(2)$因为
\[
\big|\bm{A}\big|\left|\bm{A}^{-1}\right|=1 
\Longrightarrow \left| \bm{A}^{-1}\right|=
\big | \bm{A}\big|^{-1}
\]

$(3)$仅证$\bm{AB}=\bm{I}_n$,因为
\[
1=\left|\bm{I}_n\right|=\left|\bm{A}\right|\left|\bm{B}\right|
\Longrightarrow \left|\bm{A}\right|\neq 0 \Longrightarrow \bm{A}^{-1}\text{存在}
\]
故
\[
\bm{B}=\bm{I}_n\bm{B}=\left(\bm{A}^{-1}\bm{A}\right)\bm{B}=\bm{A}^{-1}\bm{I}_n=\bm{A}^{-1}\qedhere
\]
\end{Proof}
\end{green}
\newpage
\section{初等变换与初等矩阵}
\subsection{行列式运算的性质}
\begin{formal}
\begin{proposition}[行列式的若干性质]\label{prop:行列式的若干性质}
下面$\bm{A}$、$\bm{B}$均为$n$阶方阵
,$c \in \mathbb{R} $
\begin{enumerate}[label=\textup{(\arabic*)}]
    \item 一般$ \det\left(\bm{A}
    + \bm{B}\right) \neq  \det\left(
        \bm{A}
    \right) +  \det\left(\bm{B}\right)$而是如\cref{cor:矩阵和的行列式}
    \item $ \det\left(c\bm{A}\right) = 
    c^n  \det\left(\bm{A}\right)$
    \item $ \det\left(\bm{AB}\right) = 
     \det\left(\bm{A}\right) \cdot
     \det\left(\bm{B}\right)$
    \item $ \det\left(\bm{A}^{\prime}\right) = 
     \det\left(\bm{A}\right)$
    \item $ \det\left(\overline{\bm{A}} 
    \right) = \overline{
         \det\left(\bm{A}\right)}$
    \item $ \det\left(\bm{A}^{-1}\right)
    = \left[ \det\left(\bm{A}
    \right)\right]^{-1}$
    \item $ \det\left(\bm{A}^*\right) = 
    det^{n - 1}\left(\bm{A}\right)$
\end{enumerate}
\end{proposition}
\end{formal}
\subsection{矩阵的初等变换}
\begin{formal}
\begin{definition}[矩阵的初等变换]\label{def:矩阵的初等变换}
    对矩阵的第一类、第二类、第三类初等
    行（列）变换为
    \begin{enumerate}[label=\textup{(\arabic*)}]
        \item 对换矩阵中两行（列）的位置
        \item 用一非零常数$c$乘以某一行（列）
        \item 将矩阵某一行（列）乘以一常数
        $c$加到另一行（列）上
    \end{enumerate}
\end{definition}
\end{formal}
\subsection{相抵/等价}
\begin{formal}
\begin{definition}[相抵]\label{def:相抵}
    若$\bm{A}$经过有限次的初等变换（行列均可）
    化为$\bm{B}$,则称$\bm{A}$与$\bm{B}$
    是等价的,或者说$\bm{A}$相抵于$\bm{B}$,
    记作
    \[
    \bm{A} \sim \bm{B}
    \]
\end{definition}
\end{formal}
\begin{formal}
\begin{theorem}[相抵标准型]\label{thm:相抵标准型}
    任一$\bm{A}_{m \times n}$一定相抵于
    它的相抵标准型$\bm{B}$
    \[
    \bm{B} = 
    \begin{pmatrix}
        \bm{I}_r & \bm{O} \\
        \bm{O} & \bm{O} 
    \end{pmatrix}
    \quad 
    \left(0 \leqslant 
    r \leqslant \min \,\left\{
        m,n
    \right\},r \in \mathbb{N} \right)
    \]
    并且其中的$r$不依赖于初等变换的选取
    而仅由$\bm{A}$唯一确定.
\end{theorem}
\begin{Proof}
    考虑对$\min\,\left\{m,n\right\}$进行归纳,若$\min\,\left\{m,n\right\}=0$则归纳结束.设$\min\,\left\{m,n\right\}<k$时成立,下面证明$\min,\left\{m,n\right\}=k$的情形.

    显然$\bm{A}=\bm{O}$时成立.下设$\bm{A}\neq\bm{O} $,不妨设$a_{ij}\neq 0.$那么将第$i$行与第$1$行对换,第$j$列与第$1$列对换,那么以下设$a_{11}\neq 0$
    \[
    \bm{A}=\begin{pmatrix}
        a_{11}&a_{12}&\cdots&a_{1n}\\
        a_{21}&a_{22}&\cdots&a_{2n}\\
        \vdots&\vdots&&\vdots\\
        a_{m1}&a_{m2}&\cdots&a_{mn}
    \end{pmatrix}\longrightarrow
    \begin{pmatrix}
        1&0&\cdots&0\\
        0&b_{22}&\cdots&b_{2n}\\
        \vdots&\vdots&&\vdots\\
        0&b_{m2}&\cdots&b_{mn}
    \end{pmatrix}
    \]于是根据归纳假设即证.
\end{Proof}
\end{formal}
\begin{formal}
\begin{definition}[阶梯点]\label{def:阶梯点}
    设$\bm{A}=\left(a_{ij}\right)_{m\times n}$,$\forall 1\leqslant i\leqslant m$\[
    k_i=\begin{cases*}
    +\infty&,若$a_{ij}=0,\forall 1\leqslant j\leqslant n$\\
    \min\,\left\{j\mid a_{ij}\neq 0\right\}&,若$\exists a_{ij}\neq 0$
    \end{cases*}
    \]若$\bm{A}$的第$i$行不全为零,则$a_{ik_i}$是第$i$行从左至右第一个非零元,称为第$i$行的阶梯点.
\end{definition}
\end{formal}
下面给出\cref{def:阶梯形矩阵}的严格定义
\begin{formal}
\begin{definition}[阶梯形矩阵]\label{def:阶梯形矩阵2}
    在\cref{def:阶梯点}基础上,$\bm{A}$是阶梯形矩阵当且仅当存在$0\leqslant r\leqslant m\,\textup{s.t.}$
    \[
    k_1<k_2<\cdots<k_r,k_{r+1}=\cdots=k_m=+\infty
    \]
\end{definition}
\end{formal}
\begin{formal}
\begin{theorem}[相抵于阶梯形]\label{thm:相抵于阶梯形}
    任一$\bm{A}_{m \times n}$
    在初等行变换下一定相抵于
    一个阶梯形矩阵（\cref{def:阶梯形矩阵}）.
\end{theorem}
\begin{Proof}
对列数$n$进行归纳,$n=0$表示归纳结束,设列数小于$n$时成立,下面证明$n$时成立.

一方面,如果$\bm{A}$第一列全为零,由归纳假设得证.另一方面,如果$\bm{A}$第一列不全为零,用行变换将非零元换到$\left(1,1\right)$位置,消去同列其他元素得到一个$\left(m-1\right)\times\left(n-1\right)$阶矩阵,由归纳假设得证.于是证毕.
\end{Proof}
\end{formal}
\subsection{初等矩阵}
\begin{formal}
\begin{definition}[初等矩阵]\label{def:初等矩阵}
    对单位阵$\bm{I}_n$分别施加第一类、第二类
    、第三类初等变换后得到
    的矩阵分别称为第一类、第二类、
    第三类初等矩阵.
\paragraph{第一类初等矩阵}
第一类初等矩阵$P_{ij}$表示将单位阵$\bm{I}_n$
的$i$行与第$j$列（或者是第$i$列与第$j$列）对换得到的
矩阵
\[
\bm{P}_{ij} = 
\begin{pmatrix}
    1 & & & & & & \\
     & \ddots & & &&&\\
     &&0&\cdots&1&&\\
    &&\vdots&&\vdots&&\\
    && 1 & \cdots & 0 &&\\
    &&&&&\ddots&\\
    &&&&&&1
\end{pmatrix}
\]
\paragraph{第二类初等矩阵}
第二类初等矩阵$P_i\left(c\right)$表示
用常数$c \in \mathbb{R}\left(c \neq 0\right)$
单位阵第$i$行（或者是第$i$列）得到的矩阵
\[
\bm{P}_i\left(c\right) = 
\begin{pmatrix}
    1 & & & & \\
      &\ddots&&&\\
    &&c&&\\
    &&&\ddots&\\
    &&&&1
\end{pmatrix}
\]
\paragraph{第三类初等矩阵}
第三类初等矩阵$\bm{T}_{ij}\left(c\right)$
表示将单位阵第$i$行（或者是第$i$列）乘以一常数$c$
后加到第$j$行（或者是第$j$列）得到的矩阵
\[
\bm{T}_{ij} \left(c\right)= 
\begin{pmatrix}
    1 & & & & & & \\
     & \ddots & & &&&\\
     &&1&\cdots&0&&\\
    &&\vdots&&\vdots&&\\
    && c & \cdots & 1 &&\\
    &&&&&\ddots&\\
    &&&&&&1
\end{pmatrix}
\]
\end{definition}
\end{formal}
\begin{formal}
\begin{theorem}[初等变换与初等矩阵]\label{thm:初等变换与初等矩阵}
    对于$m \times n$矩阵$\bm{A}$
    ,对$\bm{A}$作一次初等行变换
    得到的矩阵等于
    用一个$m$阶对应的初等
    矩阵左乘$\bm{A}$的积；
    对$\bm{A}$作一次初等列变换
    得到的矩阵等于
    用一个$m$阶对应的初等
    矩阵右乘$\bm{A}$的积.
\end{theorem}
\end{formal}
\begin{red}
\begin{remark}
    初等行（列）变换等价于
左（右）乘对应的初等矩阵.
\end{remark}
\end{red}
\begin{green}
\begin{corollary}[初等矩阵的逆阵]\label{cor:初等矩阵的逆阵}
    初等矩阵均为可逆阵
    且其逆阵仍为同类初等矩阵
    \[
    \bm{P}_{ij}^{-1} = \bm{P}_{ij} \quad
    \bm{P}_i\left(c\right)^{-1}
    = \bm{P}_i\left(\frac{1}{c}\right)
    \quad
    \bm{T}_{ij}\left(c\right)^{-1}
    = \bm{T}_{ij}\left(-c\right)
    \]
\end{corollary}
\end{green}
\begin{green}
\begin{corollary}[初等矩阵的行列式的值]\label{cor:初等矩阵的行列式的值}
    \[
    \left|\bm{P}_{ij}\right|=-1,
    \left|\bm{P}_i\left(c\right)\right|=c,
    \left|\bm{T}_{ij}\left(c\right)\right|=1
    \]
\end{corollary}
\end{green}
\begin{green}
\begin{corollary}[第三类初等变换]\label{cor:第三类初等变换}
    由前文\cref{thm:初等变换与初等矩阵},\cref{cor:初等矩阵的行列式的值}和\cref{thm:方阵积的行列式}可知对矩阵施加第三类初等变换不改变其行列式的值.
\end{corollary}
\end{green}
\begin{green}
\begin{corollary}[初等变换不改变矩阵的奇异性]\label{cor:初等变换不改变矩阵的奇异性}
    由\cref{cor:第三类初等变换}知初等变换不改变矩阵的奇异与否.
\end{corollary}
\end{green}
\subsection{相抵}
在说明相抵这一关系的性质前,有必要对
“关系”作一个简单的说明.
\begin{formal}
    \begin{definition}[关系]\label{def:关系}
考虑集合$A$自身的笛卡尔积
\[
A \times A = \left\{
    \left(a,b\right) | 
    a,b \in A
\right\}
\]
给定$R \subseteq  A \times A$
.若$\left(a,b\right) \in R$
,称$a$,$b$有关系$R$,记作
\[
a R b
\]
    \end{definition}\end{formal}
    \begin{formal}
    \begin{definition}[等价关系]\label{def:等价关系}
接下来定义满足下列性质的关系为等价关系：
\begin{enumerate}[label=\textup{(\arabic*)}]
    \item 自反性： $aRa$
    \item 对称性：$aRb\Longrightarrow 
    bRa$
    \item 传递性：若$aRb$且
    $bRc$,则$aRc$ 
\end{enumerate}
    \end{definition}\end{formal}
    \begin{formal}
    \begin{proposition}[相抵的性质]\label{prop:相抵的性质}
相抵是一种等价关系,即满足上述三条性质
 \begin{enumerate}[label=\textup{(\arabic*)}]
   \item 自反性：$ \bm{A} 
    \sim \bm{A}$
     \item 对称性：$\bm{A} \sim
     \bm{B}\Longrightarrow 
     \bm{B}
    \sim \bm{A}$
     \item 传递性：若$\bm{A}\sim
     \bm{B}$且$
     \bm{B} \sim \bm{C}
     $,则$\bm{A} \sim
    \bm{C}$ 
\end{enumerate}
 \end{proposition}
    \end{formal}
    \begin{formal}
    \begin{proposition}[方阵的若干等价结论]\label{prop:方阵的若干等价结论}
如果$\bm{A}$为$n$阶方阵,则下列是等价的
\begin{enumerate}[label=\textup{(\arabic*)}]
    \item $\bm{A}$是非异阵
    \item $\bm{A}$ 的相抵标准型为$\bm{I}_n$
    \item $\bm{A}$仅通过初等行
    （列）变换就可以化为$\bm{I}_n$
    \item $\bm{A}$等于若干个初等矩阵的积
\end{enumerate}
    \end{proposition}
    \begin{Proof}
        $(1)\Longrightarrow (2)$设$\bm{A}$的相抵标准型为$\begin{pmatrix}
            \bm{I}_r & \bm{O} \\
            \bm{O} & \bm{O}
        \end{pmatrix}$,由$\bm{A}$非异及\cref{cor:初等变换不改变矩阵的奇异性}知该相抵标准型非奇异,故$r=n$.

        $(2)\Longrightarrow (3)$设存在初等阵$\bm{P}_1,\bm{P}_2,\cdots,\bm{P}_r,\bm{Q}_1,\bm{Q}_2,\cdots,\bm{Q}_s\,\textup{s.t.}$
        \[
        \bm{P}_1\bm{P}_2\cdots\bm{P}_r\bm{A}\bm{Q}_1\bm{Q}_2\cdots\bm{Q}_s=\bm{I}_n
        \]作逆阵即证.

        $(3)\Longrightarrow (4)$依据上一证明,显然.

        $(4)\Longrightarrow (1)$由\cref{cor:初等矩阵的逆阵}知初等矩阵均非异,故$\bm{A}$非异.
    \end{Proof}
    \end{formal}
    \begin{green}
    \begin{corollary}[相抵的等价表述]\label{cor:相抵的等价表述}
        对于$\bm{A}=\left(a_{ij}\right)_{m\times n}\in M_{m\times n}\left(\mathbb{R}\right)$,若存在初等阵$\bm{P}_1,\bm{P}_2,\cdots,\bm{P}_r,\bm{Q}_1,\bm{Q}_2,\cdots,\bm{Q}_s\,\textup{s.t.}$
        \[
        \bm{P}_1\bm{P}_2\cdots\bm{P}_r\bm{A}\bm{Q}_1\bm{Q}_2\cdots\bm{Q}_s=\bm{B}
        \]则称$\bm{A}$与$\bm{B}$相抵.
    \end{corollary}
    \begin{Proof}
    根据\cref{thm:初等变换与初等矩阵}和\cref{thm:相抵标准型}易证.
    \end{Proof}
    \end{green}
\newpage
\section{矩阵积的行列式与初等变换求逆阵}
\subsection{$n$阶方阵的积的行列式}
\begin{formal}
\begin{theorem}[方阵积的行列式]\label{thm:方阵积的行列式2}
    设$\bm{A}$、$\bm{B}$均为$n$阶方阵
    ,则
    \[
    \left|\bm{A}\bm{B}\right| = 
    \left|\bm{A}\right|\left|\bm{B}\right|
    \]
\end{theorem}
\end{formal}
\subsection{初等变换求逆阵}
\begin{green}
\begin{lemma}[相抵标准型]\label{lem:相抵标准型}
    对于$n$阶方阵$\bm{A}$（不一定可逆）,总存在$n$
    阶非异阵$\bm{P}$、$\bm{Q}$使得
    \[
    \bm{PAQ} = \begin{pmatrix}
        \bm{I}_r & \bm{O} \\
        \bm{O} & \bm{O}
    \end{pmatrix}
    \]
\end{lemma}
\end{green}
不妨将这个问题用于解方程
\[
\bm{A}\bm{X} = \bm{B}
\]
显然有
\[
\bm{X} = \bm{A}^{-1}\bm{B}
\]
只需构造
\[\qquad\quad
\left (
    \begin{array}{c:c}
        \begin{matrix}
            \bm{A}
        \end{matrix}
        &
        \begin{matrix}
            \bm{B}
        \end{matrix}
    \end{array}
\right )_{n \times \left(
    n + m
\right)}
\]
用初等行变换将$\bm{A}$化为
$\bm{I}_n$时右边得到
$\bm{X} = \bm{A}^{-1}\bm{B}$

特别地,求逆阵时构造
\[\qquad\quad
\left (
    \begin{array}{c:c}
        \begin{matrix}
            \bm{A}
        \end{matrix}
        &
        \begin{matrix}
            \bm{I}_n
        \end{matrix}
    \end{array}
\right )_{n \times \left(
    n + m
\right)}
\]
另外,欲使用初等列变换,应构造
\[
\left ( 
    \begin{array}{c}
        \bm{A} \\
        \hdashline
        \bm{B}
    \end{array}
\right )
\]
形式.
\begin{red}
\begin{remark}
    该方法本质上就是初等变换.
\end{remark}
\end{red}
\newpage
\section{分块矩阵}
\subsection{分块矩阵}
\begin{formal}
\begin{definition}[分块矩阵]\label{def:分块矩阵}
    将矩阵分割为几块,得到一个分块矩阵
    \[
    \bm{A} = \left(\bm{A}_{ij}\right)_
    {r \times s} = 
    \begin{pmatrix}
        \bm{A}_{11} & \bm{A}_{12}&\cdots&\bm{A}_{1s}\\
        \bm{A}_{21} & \bm{A}_{22}&\cdots&\bm{A}_{2s}\\
        \vdots & \vdots & & \vdots\\
        \bm{A}_{r1} & \bm{A}_{r2}&\cdots&\bm{A}_{rs}
    \end{pmatrix}
    \]
其中\[
\bm{A}_{i1},\bm{A}_{i2},\cdots,\bm{A}_{is} 
\quad \left(i = 1,2,\cdots,r\right)
\]
为第$i$分块行.分块$\bm{A}_{ij}=\left(a_{lk}\right)_{m_i\times n_j},\forall 1\leqslant i\leqslant r,1\leqslant j\leqslant n$,其中
\[
\left(m_1,m_2,\cdots,m_r\right),\left(
    n_1,n_2,\cdots,n_s
\right)
\]
分别是$\bm{A}$的行、列分块方式.
\end{definition}
\end{formal}
\begin{formal}
\begin{definition}[分块矩阵的相等]\label{def:分块矩阵的相等}
    设$\bm{A} = \left(\bm{A}_{ij}\right)_{
    r \times s
}$、$\bm{B} = 
\left(\bm{B}_{ij}\right)_{
    k \times l
}$,二者相等当且仅当
\begin{equation*}
    \begin{cases*}
        r  = k,s = l  \\
        \bm{A}_{ij} =
        \bm{B}_{ij}\quad
        \left(\forall 1 \leqslant 
        i \leqslant r,1 \leqslant  j
        \leqslant s\right)
    \end{cases*}
\end{equation*}
\end{definition}
\end{formal}
\begin{red}
\begin{remark}
    分块矩阵实际上只是一种方便的记号,并没有实质上的变化.
\end{remark}
\end{red}
\subsection{运算}
\begin{formal}
\begin{definition}[运算]\label{def:运算}
定义分块矩阵相关运算为：
\begin{enumerate}[label=\textup{(\arabic*)}]
    \item 加法：若$\bm{A} = \left(\bm{A}_{ij}\right)_{
    r \times s
}$、$\bm{B} = 
\left(\bm{B}_{ij}\right)_{
    r \times s
}$行、列分块方式相同
,定义加法
\[
\bm{A}+\bm{B} = \left(
    \bm{A}_{ij} + \bm{B}_{ij}
\right)_{r \times s}
\]
\item 数乘：$\bm{A} = \left(\bm{A}_{ij}\right)_{
    r \times s
}$,$c \in \mathbb{R} $
,数乘
\[
c\bm{A} = \left(c\bm{A}_{ij}\right)
_{r \times s}
\]
\item 乘法：设$\bm{A} = \left(\bm{A}_{ij}\right)_{
    r \times s
}$、$\bm{B} = 
\left(\bm{B}_{ij}\right)_{
    s \times k
}$,且$\bm{A}$的列分块方式与
$\bm{B}$的行分块方式相同,那么
\[
    \bm{AB} = \left(\bm{C}_{ij}
    \right)_{r \times k} 
    = \left(
        \sum_{l = 1}^{s}\bm{A}_{il}
        \bm{B}_{lj}
    \right)_{
        r \times k
    }
\]
特别地,如果有\cref{def:形式行向量}一样的形式记号即具有行分块
$\bm{A} = \left(a_{ij}\right)_{
    m \times n
} = \left(
    \bm{\alpha }_1,\cdots
    ,\bm{\alpha }_m
\right)^{\prime}$,列分块
$\bm{B} = \left(b_{ij}\right)_{
    n \times p
} = \left(
    \bm{\beta  }_1,\cdots
    ,\bm{\beta  }_p
\right)$,那么
\[
\bm{AB} = \left(\bm{\alpha }_1
\bm{B},\cdots,
\bm{\alpha }_m\bm{B}\right)^{\prime} =
\left(
    \bm{A}\bm{\beta }_1,\cdots
    ,\bm{A}\bm{\beta }_p
\right)
\]
\item 转置：设$\bm{A} = \left(\bm{A}_{ij}\right)_{
    r \times s
}$,转置为
\[
\bm{A}^{\prime} = \left(
    \bm{A}_{ij}
\right)_{s \times r}
\]
\item 共轭：设复矩阵$\bm{A} = \left(\bm{A}_{ij}\right)_{
    r \times s
}$,共轭为
\[
\overline{\bm{A}} = 
\left(\overline{\bm{A}_{ij}} \right)
_{r \times s}
\]   
\end{enumerate}
\end{definition}
\end{formal}
\begin{brown}
\begin{example}
    若$\bm{A} = \mathrm{diag}\left\{
        \bm{A}_1,\cdots,\bm{A}_k
    \right\}$,$
    \bm{B}= \mathrm{diag}\left\{
        \bm{B}_1,\cdots,\bm{B}_k
    \right\}$
    ,则
    \[
    \bm{AB} = \mathrm{diag}\left\{
        \bm{A}_1\bm{B}_1,\cdots
        ,\bm{A}_k\bm{B}_k
    \right\}
    \]
\end{example}
\end{brown}\begin{green}
\begin{corollary}[分块矩阵的逆阵]\label{cor:分块矩阵的逆阵}
    若$\bm{A} = \mathrm{diag}\left\{
        \bm{A}_1,\cdots,\bm{A}_k
    \right\}$可逆则
    \[
    \bm{A}^{-1} = 
    \mathrm{diag}\left\{
        \bm{A}_1^{-1},\cdots
        ,\bm{A}_k^{-1}
    \right\}
    \]
\end{corollary}
\end{green}
\subsection{分块初等矩阵}
\begin{formal}
\begin{definition}[分块单位阵]\label{def:分块单位阵}
    分块单位阵为
\[
\bm{I} = \mathrm{diag}\left\{
    \bm{I}_{m_1},\cdots,
    \bm{I}_{m_n}
\right\}
\]
\end{definition}
\end{formal}

\begin{formal}
\begin{definition}[分块初等矩阵]\label{def:分块初等矩阵}
    三类定义与初等矩阵类似.事
    实上,分块只是一种写法,并
    不是什么新奇东西.只是操作的一行或一列变成了几行或几列.
\paragraph{第一类分块初等矩阵}对换$\bm{I}$的第$i$分块行（列）与第$j$分块行（列）得到.
\[
\bm{P}_{ij} = 
\begin{pmatrix}
    \bm{I}_{m_{1}} & & & & & & \\
     & \ddots & & &&&\\
     &&0&\cdots&\bm{I}_{m_j}&&\\
    &&\vdots&&\vdots&&\\
    && \bm{I}_{m_i} & \cdots & 0 &&\\
    &&&&&\ddots&\\
    &&&&&&\bm{I}_{m_n}
\end{pmatrix}
\]
\paragraph{第二类分块初等矩阵}
$ \det\left(\bm{M}\right) \neq 0$下以$\bm{M}$左（右）乘
$\bm{I}$的第$i$分块行（列）得到.
\[
\bm{P}_i\left(\bm{M}\right) = 
\begin{pmatrix}
    \bm{I}_{m_1} & & & & \\
      &\ddots&&&\\
    &&\bm{M}&&\\
    &&&\ddots&\\
    &&&&\bm{I}_{m_n}
\end{pmatrix}
\]
\newpage
\paragraph{第三类分块初等矩阵}
以矩阵$\bm{M}$左（右）乘$\bm{I}$的第$i$分块行（列）后加到第
$j$分块行（列）得到.
\[
\bm{T}_{ij} \left(\bm{M}\right)= 
\begin{pmatrix}
    \bm{I}_{m_1} & & & & & & \\
     & \ddots & & &&&\\
     &&\bm{I}_{m_i}&\cdots&0&&\\
    &&\vdots&&\vdots&&\\
    && \bm{M} & \cdots & \bm{I}_{m_j} &&\\
    &&&&&\ddots&\\
    &&&&&&\bm{I}_{m_n}
\end{pmatrix}
\]
\end{definition}
\end{formal}
\begin{formal}
\begin{theorem}[分块初等矩阵的逆阵]\label{thm:分块初等矩阵的逆阵}
分块初等矩阵均可逆且逆阵为同类分块初等矩阵
\[
\bm{P}_{ij}^{-1} = \bm{P}_{ij}^{\prime} \quad
\bm{P}_i\left(\bm{M}\right)^{-1}
=\bm{P}_i\left(\bm{M}^{-1}\right) \quad
\bm{T}_{ij}\left(\bm{M}\right)^{-1}
= \bm{T}_{ij}\left(-\bm{M}\right)
\]
\end{theorem}
\end{formal}
\begin{formal}
\begin{theorem}[分块初等矩阵的行列式]\label{thm:分块初等矩阵的行列式}
\[
 \det\left(\bm{P}_{ij}\right) = \left(-1
\right)^l, \det\left(\bm{P}_i\left(\bm{M}\right)\right)
=  \det\left(\bm{M}\right),
 \det\left(\bm{T}_{ij}\left(\bm{M}
\right)\right) = 1
\]
其中
\[
l = m_i m_j + \left(m_i + m_j\right)
\sum_{i < r <j}m_r
\]
\end{theorem}
\end{formal}
\subsection{分块初等变换}
\begin{formal}\begin{theorem}[分块初等变换与分块初等矩阵]\label{thm:分块初等变换与分块初等矩阵}
    类似于\cref{thm:初等变换与初等矩阵},每一种分块初等行（列）变换等价于左（右）
乘这种变换对应的分块初等矩阵.
\end{theorem}\end{formal}
\begin{green}
\begin{corollary}[分块初等变换与行列式]\label{cor:分块初等变换与行列式}
    第三类分块初等变换
不会改变矩阵行列式的值.
\end{corollary}
\end{green}
\subsection{行列式的降阶公式}
\begin{green}
\begin{lemma}[分块三角行列式]\label{lem:分块三角行列式}
    对分块三角行列式有
    \[
    \left|\bm{G}\right| = 
    \begin{vmatrix}
        \bm{A} & \bm{B} \\
        \bm{O} & \bm{C} 
    \end{vmatrix} 
    = \left|\bm{A}\right|
    \left|\bm{C}\right|
    \]
    \[
    \left|\bm{H}\right| = 
    \begin{vmatrix}
        \bm{A} & \bm{O} \\
        \bm{B} & \bm{C} 
    \end{vmatrix} 
    = \left|\bm{A}\right|
    \left|\bm{C}\right|
    \]
\end{lemma}
\begin{Proof}
    即\cref{cor:分块行列式},
利用Laplace定理是显然的.
\end{Proof}
\end{green}
\begin{formal}
\begin{theorem}[分块矩阵行列式降阶公式]\label{thm:分块矩阵行列式降阶公式}
    对于矩阵
    \[
    \bm{M} = \begin{pmatrix}
        \bm{A} & \bm{B} \\
        \bm{C} & \bm{D} 
    \end{pmatrix}
    \]
$1.$若$\bm{A}$可逆,则
\[
\left|\bm{M}\right| = \left|\bm{A}
\right|\left|\bm{D} - \bm{C}\bm{A}
^{-1}\bm{B}\right|
\]
$2.$若$\bm{D}$可逆,则\[
\left|\bm{M}\right| = \left|
    \bm{D}
\right|\left|
    \bm{A} - \bm{B}\bm{D}^{-1}
    \bm{C}
\right|
\]
$3.$若$\bm{A}$、$\bm{D}$均可逆,则
\[
\left|\bm{A}
\right|\left|\bm{D} - \bm{C}\bm{A}
^{-1}\bm{B}\right|
=
\left|
    \bm{D}
\right|\left|
    \bm{A} - \bm{B}\bm{D}^{-1}
    \bm{C}
\right|
\]
\end{theorem}
\begin{Proof}
证明是比较显然的,构造\cref{lem:分块三角行列式}形式即可.
\end{Proof}
\end{formal}
\begin{brown}
\begin{example}
    计算下面矩阵的行列式
    \[
    \bm{M}=\begin{pmatrix}
        a_1^2&a_1a_2+1&\cdots&a_1a_n+1\\
        a_2a_1+1&a_2^2&\cdots&a_2a_n+1\\
        \vdots&\vdots&&\vdots\\
        a_na_1+1&a_na_2+1&\cdots&a_n^2
    \end{pmatrix}
    \]
    \begin{solution}
    考虑
    \begin{align*}
    \bm{M}&=\begin{pmatrix}
        a_1^2+1&a_1a_2+1&\cdots&a_1a_n+1\\
        a_2a_1+1&a_2^2+1&\cdots&a_2a_n+1\\
        \vdots&\vdots&&\vdots\\
        a_na_1+1&a_na_2+1&\cdots&a_n^2+1
    \end{pmatrix}-\bm{I}_n\\
    &=-\bm{I}_n+\begin{pmatrix}
        a_1&1\\
        a_2&1\\
        \vdots&\vdots\\
        a_n&1
    \end{pmatrix}\begin{pmatrix}
        a_1&a_2&\cdots&a_n\\
        1&1&\cdots&1
    \end{pmatrix}\\
    &=-\bm{I}_n+\begin{pmatrix}
        a_1&1\\
        a_2&1\\
        \vdots&\vdots\\
        a_n&1
    \end{pmatrix}
    \bm{I}_2^{-1}\begin{pmatrix}
        a_1&a_2&\cdots&a_n\\
        1&1&\cdots&1
    \end{pmatrix}
\end{align*}
由降阶公式有
\begin{align*}
    \left|\bm{M}\right|&=
    \left|\bm{I}_2\right|^{-1}\cdot 
    \left|-\bm{I}_n\right|\cdot
    \left|
        \bm{I}_2+\begin{pmatrix}
        a_1&a_2&\cdots&a_n\\
        1&1&\cdots&1
    \end{pmatrix}\left(-\bm{I}_n\right)^{-1}
    \begin{pmatrix}
        a_1&1\\
        a_2&1\\
        \vdots&\vdots\\
        a_n&1
    \end{pmatrix}
    \right|\\
    &=\left(-1\right)^n\left|
        \bm{I}_2-\begin{pmatrix}
            \sum\limits_{i=1}^{n}a_i^2&\sum\limits_{i=1}^{n}a_i\\
            \sum\limits_{i=1}^{n}a_i&n
        \end{pmatrix}
    \right|\\
    &=\left(-1\right)^n\left(
        \left(1-n\right)
        \left(1-\sum\limits_{i=1}^{n}a_i^2\right)-\left(
            \sum\limits_{i=1}^{n}a_i
        \right)^2
    \right)
\end{align*}
    \end{solution}
\end{example}
\end{brown}
\begin{brown}
    \begin{example}
        $\bm{A}$、$\bm{D}$均可逆,
        求分块矩阵的逆阵
        \[
        \begin{pmatrix}
            \bm{A}&\bm{B}\\
            \bm{O}&\bm{D}
        \end{pmatrix}
        \]
        \newpage
        \begin{solution}
        设$\bm{A}$、$\bm{D}$的阶数分别为$m$、$n$,那么
        \begin{align*}
        \left (
        \begin{array}{c:c}
            \begin{matrix}
                \bm{A}&\bm{B}\\
                \bm{O}&\bm{D}
            \end{matrix}
            &
            \begin{matrix}
                \bm{I}_m&\bm{O}\\
                \bm{O}&\bm{I}_n
            \end{matrix}
        \end{array}
        \right )
        &\longrightarrow
        \left (
        \begin{array}{c:c}
            \begin{matrix}
                \bm{A}&\bm{O}\\
                \bm{O}&\bm{D}
            \end{matrix}
            &
            \begin{matrix}
                \bm{I}_m&-\bm{BD}^{-1}\\
                \bm{O}&\bm{I}_n
            \end{matrix}
        \end{array}
        \right )\\
        &\longrightarrow
        \left (
        \begin{array}{c:c}
            \begin{matrix}
                \bm{I}_m&\bm{O}\\
                \bm{O}&\bm{I}_n
            \end{matrix}
            &
            \begin{matrix}
                \bm{A}^{-1}&-\bm{A}^{-1}\bm{BD}^{-1}\\
                \bm{O}&\bm{D}^{-1}
            \end{matrix}
        \end{array}
        \right )
    \end{align*}
        \end{solution}
    \end{example}
\end{brown}
\begin{brown}
\begin{example}
    设$n$阶方阵$\bm{A}$、$\bm{B}$,求证：
    \[
    \begin{vmatrix}
        \bm{A}&\bm{B}\\
        \bm{B}&\bm{A}
    \end{vmatrix}
    =
    \left|\bm{A}+\bm{B}\right|
    \left|\bm{A}-\bm{B}\right|
    \]
    \begin{solution}
        作第三类分块初等变换
        \begin{align*}
            \begin{vmatrix}
        \bm{A}&\bm{B}\\
        \bm{B}&\bm{A}
    \end{vmatrix}
    &=
    \begin{vmatrix}
        \bm{A}+\bm{B}&\bm{A}+\bm{B}\\
        \bm{B}&\bm{A}
    \end{vmatrix}\\
    &=\begin{vmatrix}
        \bm{A}+\bm{B}&\bm{O}\\
        \bm{B}&\bm{A}-\bm{B}
    \end{vmatrix}\\
    &=
    \left|\bm{A}+\bm{B}\right|
    \left|\bm{A}-\bm{B}\right|
        \end{align*}
    \end{solution}
\end{example}
\end{brown}
\begin{brown}
\begin{example}\label{ex:1}
    设$n$阶复矩阵$\bm{A}$、$\bm{B}$,求证：
    \[
    \begin{vmatrix}
        \bm{A}&-\bm{B}\\
        \bm{B}&\bm{A}
    \end{vmatrix}=\left|
        \bm{A}+\mathrm{i}\bm{B}
    \right|\left|\bm{A}-\mathrm{i}\bm{B}\right|
    \]
    \begin{solution}
    \begin{align*}
            \begin{vmatrix}
        \bm{A}&-\bm{B}\\
        \bm{B}&\bm{A}
    \end{vmatrix}
    &=
    \begin{vmatrix}
        \bm{A}+\mathrm{i}\bm{B}&\mathrm{i}\bm{A}-\bm{B}\\
        \bm{B}&\bm{A}
    \end{vmatrix}\\
    &=\begin{vmatrix}
        \bm{A}+\mathrm{i}\bm{B}&\bm{O}\\
        \bm{B}&\bm{A}-\mathrm{i}\bm{B}
    \end{vmatrix}\\
    &=
    \left|\bm{A}+\mathrm{i}\bm{B}\right|
    \left|\bm{A}-\mathrm{i}\bm{B}\right|
        \end{align*}
    \end{solution}
\end{example}
\end{brown}
\begin{brown}
    \begin{example}
        设$n$阶方阵$\bm{A}$、$\bm{B}$且$\bm{AB}=\bm{BA}$,求证：
        \[
        \begin{vmatrix}
            \bm{A}&-\bm{B}\\
            \bm{B}&\bm{A}
        \end{vmatrix}=\left|\bm{A}^2+\bm{B}^ 2\right|
        \]
    \end{example}
    \begin{Proof}
    由\cref{ex:1}立得.
    \end{Proof}
\end{brown}
\newpage
\section{Cauchy-Binet公式}
\begin{red}
\begin{remark}
    \textup{Cauchy-Binet}公式研究的是非方阵的矩阵积的行列式.
\end{remark}
\end{red}
\subsection{n阶方阵的积的行列式}
\begin{formal}
\begin{theorem}[方阵积的行列式]
    对于$n$阶方阵$\bm{A}$、$\bm{B}$
    \[
    \left|\bm{AB}\right|
    = 
    \left|\bm{A}\right|
    \left|\bm{B}\right|
    \]
\end{theorem}
\end{formal}
\subsection{Cauchy-Binet公式}
\setcounter{equation}{0}%重置计数器,方程下计数器为equation
\begin{formal}
\begin{theorem}[Cauchy-Binet公式]\label{thm:Cauchy-Binet公式}
    对于矩阵$\bm{A}_{m \times n}$、$
    \bm{B}_{n \times m}$
    \begin{enumerate}[label=\textup{(\arabic*)}]
        \item $m > n,\left|\bm{AB}\right| = 0;$
        \item $\displaystyle
        m \leqslant n,\left|\bm{AB}\right|
        = \sum_{1\leqslant j_1 < j_2<\cdots
        < j_m \leqslant n}
        \bm{A}\begin{pmatrix}
            1 & 2&\cdots & m \\
            j_1 & j_2&\cdots & j_m
        \end{pmatrix}
        \bm{B}\begin{pmatrix}
            j_1 & j_2&\cdots & j_m \\
            1 & 2&\cdots & m 
        \end{pmatrix}$
    \end{enumerate}
\end{theorem}
\begin{Proof}
构造如下$m+n$阶方阵
\[
\begin{bmatrix}
    \bm{A}&\bm{O}\\
    -\bm{I}_n&\bm{B}
\end{bmatrix}
\]

主要思路仍然是第三类分块初等变换
\begin{align*}
    \begin{vmatrix}
    \bm{A}&\bm{O}\\
    -\bm{I}_n&\bm{B}
\end{vmatrix}=\begin{vmatrix}
    \bm{O}&\bm{AB}\\
    -\bm{I}_n&\bm{B}
\end{vmatrix}
\end{align*}
均对前$m$行作\cref{thm:Laplace定理}Laplace展开
\begin{align*}
    RHS&=\left|\bm{AB}\right|\left(-1\right)^{
        1+2+\cdots+m+\left(n+1\right)+\cdots+\left(n+m\right)
    }\left|-\bm{I}_n\right|\\
    &=\left(-1\right)^{n\left(m+1\right)}\left|\bm{AB}\right|
\end{align*}

当$m>n$时,则包含在前$m$行的任一$m$阶子式
至少有一列全为0即行列式值为0.由\cref{thm:Laplace定理}Laplace定理知
\[
LHS = 0
\]
故$\left|\bm{AB}\right|=0$.

当$m\leqslant n$时,由\cref{thm:Laplace定理}Laplace定理（代数余子式记作$\bm{C}$）
\[
LHS =\sum_{1\leqslant j_1<j_2<\cdots<j_m\leqslant n}\bm{A}\begin{pmatrix}
    1&2&\cdots&m\\
    j_1&j_2&\cdots&j_m
\end{pmatrix}\bm{C}\begin{pmatrix}
    1&2&\cdots&m\\
    j_1&j_2&\cdots&j_m
\end{pmatrix}
\]
又代数余子式
\begin{align*}
    \bm{C}\begin{pmatrix}
    1&2&\cdots&m\\
    j_1&j_2&\cdots&j_m
\end{pmatrix}=\left(-1
\right)^{1+2+\cdots+m+j_1
+j_2+\cdots+j_m}\left|-\bm{e}_{i_1},-\bm{e}_{i_2},\cdots,-\bm{e}_{i_{n-m}},\bm{B}\right|
\end{align*}
其中,$1\leqslant i_1
<i_2<\cdots<i_{n-m}
\leqslant n$为
$1\leqslant j_1
<j_2<\cdots<j_m
\leqslant n$的余指标.$
\bm{e}_i$为对应的标准单位
列向量.

记\[
\bm{N}=\left(-\bm{e}_{i_1},-\bm{e}_{i_2},\cdots,-\bm{e}_{i_{n-m}},\bm{B}\right)
\]
对前$n-m$列作\cref{thm:Laplace定理}Laplace展开
\begin{align*}
    \left|\bm{N}\right|=\left|-\bm{I}_{n-m}\right|
    \left(-1\right)^{i_1+i_2+\cdots+i_{n-m}+1+2+\cdots+\left(n-m\right)}\bm{B}\begin{pmatrix}
        j_1&j_2&\cdots&j_m\\
        1&2&\cdots&m
    \end{pmatrix}
\end{align*}
于是有
\[
LHS = \sum_{1\leqslant j_1<j_2<\cdots<j_m\leqslant n}
\left(-1\right)^l\bm{A}\begin{pmatrix}
    1&2&\cdots&m\\
    j_1&j_2&\cdots&j_m
\end{pmatrix}\bm{B}\begin{pmatrix}
    j_1&j_2&\cdots&j_m\\
    1&2&\cdots&m
\end{pmatrix}
\]
其中$l$为前面一系列指数和.
只需证明$n\left(m+1\right)$
和$l$模2同余或者二者之和为偶数.
因为
\begin{align*}
    l+n\left(m+1\right)&=
    1+2+\cdots+m\\
    &+i_1+i_2+\cdots+i_{n-m}+j_1+j_2+\cdots+j_m\\
    &+1+2+\cdots+\left(n-m\right)+\left(n-m\right)
    \\
    &+n\left(m+1\right)
\end{align*}
因为$1\leqslant i_1
<i_2<\cdots<i_{n-m}
\leqslant n$为
$1\leqslant j_1
<j_2<\cdots<j_m
\leqslant n$的余指标,故
\begin{align*}
    l+n\left(m+1\right)&=
    1+2+\cdots+m\\
    &+1+2+\cdots+n\\
    &+1+2+\cdots+\left(n-m\right)+\left(n-m\right)
    \\
    &+n\left(m+1\right)\\
    &\equiv \left(m+1\right)+\cdots+n+1+2+\cdots+\left(n-m-1\right)+n\left(m+1\right)\pmod{2}\\
    &\equiv n\left(m+1\right)+
    n\left(n-m\right)
    \pmod{2}\\
    &\equiv n\left(n+1\right)\pmod{2}\equiv 0
\end{align*}
综上,Cauchy-Binet公式证毕.
\end{Proof}
\end{formal}
可以推广求得$\bm{AB}$的$r$阶子式
\setcounter{equation}{0}%重置计数器,方程下计数器为equation
\begin{green}
\begin{corollary}[Cauchy-Binet公式推广到子式]\label{cor:Cauchy-Binet公式推广到子式}
    对于矩阵$\bm{A}_{m \times n}$、$
    \bm{B}_{n \times m}$,它们的积的$r\left(
        1 \leqslant  
    r \leqslant m\right)$阶子式
    \begin{enumerate}[label=\textup{(\arabic*)}]
        \item $r > n$,则$\bm{AB}$
            任一
        $r$阶子式为零；
        \item $ r \leqslant n$,则$\bm{AB}$
            的
        $r$阶子式
        \begin{align*}
        &\bm{AB}\begin{pmatrix}
            i_1 &i_2& \cdots & i_r \\
            j_1 &j_2& \cdots & j_r
        \end{pmatrix}\\
        =&
        \sum_{1 \leqslant  k_1 <k_2
        < \cdots < k_r \leqslant 
        n}\bm{A}
        \begin{pmatrix}
            i_1 & i_2&\cdots & i_r \\
            k_1 & k_2&\cdots & k_r
        \end{pmatrix}
        \bm{B}\begin{pmatrix}
            k_1 & k_2&\cdots & k_r \\
            j_1 & j_2&\cdots & j_r
        \end{pmatrix}
    \end{align*}
    \end{enumerate}
\end{corollary}
\begin{Proof}
设$\bm{C}_{m}=\bm{AB}$,那么
\begin{align*}
    \bm{C}\begin{pmatrix}
        i_1&i_2&\cdots&i_r\\
        j_1&j_2&\cdots&j_r
    \end{pmatrix}=\begin{vmatrix}
        \begin{bmatrix}
            a_{i_11}&a_{i_12}&\cdots&a_{i_1n}\\
            a_{i_21}&a_{i_22}&\cdots&a_{i_2n}\\
            \vdots&\vdots&&\vdots\\
            a_{i_r1}&a_{i_r2}&\cdots&a_{i_rn}
        \end{bmatrix}
        \begin{bmatrix}
            b_{1j_1}&b_{1j_2}&\cdots&b_{1j_r}\\
            b_{2j_1}&b_{2j_2}&\cdots&b_{2j_r}\\
            \vdots&\vdots&&\vdots\\
            b_{nj_1}&b_{nj_2}&\cdots&b_{nj_r}
        \end{bmatrix}
    \end{vmatrix}
\end{align*}
由\cref{thm:Cauchy-Binet公式}Cauchy-Binet公式知：

$(1)r > n$时
\[
\bm{C}\begin{pmatrix}
        i_1&i_2&\cdots&i_r\\
        j_1&j_2&\cdots&j_r
    \end{pmatrix}=0
\]

$(2)r\leqslant n$时
\[
\bm{C}\begin{pmatrix}
        i_1&i_2&\cdots&i_r\\
        j_1&j_2&\cdots&j_r
    \end{pmatrix}=
    \sum_{1\leqslant k_1<k_2<\cdots<k_r\leqslant n}\bm{A}\begin{pmatrix}
        i_1&i_2&\cdots&i_r\\
        k_1&k_2&\cdots&k_r
    \end{pmatrix}\bm{B}\begin{pmatrix}
        k_1&k_2&\cdots&k_r\\
        j_1&j_2&\cdots&j_r
    \end{pmatrix}
\]
证毕.
\end{Proof}
\end{green}
\subsection{r阶主子式}
\begin{formal}
\begin{definition}[主子式]\label{def:主子式}
    如果
    \[
    \bm{A}\begin{pmatrix}
        i_1 & \cdots & i_r \\
        j_1 & \cdots & j_r 
    \end{pmatrix}
    \]
    有
    \[
    i_k = j_k \qquad
    \left(k = 1,\cdots ,r\right)
    \]
    称为$\bm{A}$的$r$阶主子式。
\end{definition}
\end{formal}
\begin{green}
\begin{corollary}[主子式非负]\label{cor:主子式非负}
    $\bm{A}$是一个$m \times n$实矩阵
    ,则$\bm{A}\bm{A}^{\prime}$的任一主子式均
    非负.
\end{corollary}
\begin{Proof}
由\cref{thm:Cauchy-Binet公式}Cauchy-Binet公式：$r \leqslant n$时
\[
\bm{AA}^{\prime}\begin{pmatrix}
    i_1&i_2&\cdots&i_r\\
    i_1&i_2&\cdots&i_r
\end{pmatrix}=
\sum_{1\leqslant k_1<k_2<\cdots<k_r\leqslant n}\bm{A}\begin{pmatrix}
    i_1&i_2&\cdots&i_r\\
    k_1&k_2&\cdots&k_r
\end{pmatrix}^2\geqslant 0
\]

$r > n$时任一$r$阶子式均为0亦满足.
\end{Proof}
\end{green}
\begin{green}
\begin{corollary}[复矩阵推广]\label{cor:复矩阵推广}
    事实上,对于$m \times n$复矩阵$\bm{A}$
    ,$\bm{A}\overline{\bm{A}}^{\prime} $
    任一主子式均非负
\end{corollary}
\end{green}
\subsection{应用}
\begin{green}
\begin{corollary}[方阵的积的伴随阵]
对于$n$阶方阵$\bm{A}$、
$\bm{B}$
\[
\left(\bm{AB}\right)^* = \bm{B}^*\bm{A}^*
\]
\end{corollary}
\begin{Proof}
只需比较$\left(i,j\right)$元,因为$\left(\bm{AB}\right)^*$的$
\left(i,j\right)$元是$\bm{AB}$的$\left(j,i\right)$
元的代数余子式
\[
\left(-1\right)^{i+j}\bm{AB}\begin{pmatrix}
    1&\cdots&\hat{j}&\cdots&n\\
    1&\cdots&\hat{i}&\cdots&n
\end{pmatrix}
\]
由\cref{thm:Cauchy-Binet公式}Cauchy-Binet公式知上式为
\[
\left(-1\right)^{i+j}\sum_{k=1}^{n}
\bm{A}\begin{pmatrix}
    1&\cdots&\hat{j}&\cdots&n\\
    1&\cdots&\hat{k}&\cdots&n
\end{pmatrix}\bm{B}\begin{pmatrix}
    1&\cdots&\hat{k}&\cdots&n\\
    1&\cdots&\hat{i}&\cdots&n
\end{pmatrix}=\sum_{k=1}^{n}\bm{A}_{jk}\bm{B}_{ki}
\]
这个式子即是$\bm{A}^*$的第$j$列乘上$\bm{B}^*$的第$i$行即
$\bm{B}^*\bm{A}^*$的$\left(i,j\right)$元.
\end{Proof}
\end{green}

\begin{brown}
\begin{example}[Lagrange恒等式]\label{ex:Lagrange恒等式}
$n \geqslant  2$时
\[
\left(\sum_{i = 1}^{n}a_i^2\right)
\left(\sum_{i = 1}^{n}b_i^2\right)
-
\left(\sum_{i = 1}^{n}a_ib_i\right)^2
=
\sum_{1 \leqslant  i < j \leqslant n}
\left(a_ib_j - a_jb_i\right)^2
\]
\end{example}
\begin{Proof}
\begin{align*}
    LHS &= \begin{vmatrix}
        \sum\limits_{i=1}^{n}a_i^2&\sum\limits_{i=1}^{n}a_ib_i\\
        \sum\limits_{i=1}^{n}a_ib_i&\sum\limits_{i=1}^{n}b^2_i
    \end{vmatrix}\\
    &=\left|\begin{pmatrix}
        a_1&a_2&\cdots&a_n\\
        b_1&b_2&\cdots&b_n
    \end{pmatrix}\begin{pmatrix}
        a_1&b_1\\
        a_2&b_2\\
        \vdots&\vdots\\
        a_n&b_n
    \end{pmatrix}\right|
    \\
    &\xlongequal{\text{Cauchy-Binet公式}}
    \sum_{1\leqslant i < j \leqslant n}
    \begin{vmatrix}
        a_i&a_j\\
        b_i&b_j
    \end{vmatrix}\begin{vmatrix}
        a_i&b_i\\
        a_j&b_j
    \end{vmatrix}\\
    &=\sum_{1\leqslant i < j\leqslant n}
    \left(a_ib_j-a_jb_i\right)^2
    \qedhere
\end{align*}
\end{Proof}
\end{brown}
同时立刻得到下面的不等式
\begin{brown}
\begin{example}[Cauchy-Schwarz不等式]\label{ex:Cauchy-Schwarz不等式}
有实数形式与复数形式
    \paragraph{实形式} $ a_i,b_i\in\mathbb{R},n\in\mathbb{N}^+$
\[
\left(\sum_{i = 1}^{n}a_i^2\right)
\left(\sum_{i = 1}^{n}b_i^2\right)
\geqslant 
\left(\sum_{i = 1}^{n}a_ib_i\right)^2
\]
\paragraph{复形式}$a_i,b_i\in \mathbb{C},n \geqslant 2$
\[
\left(\sum_{i = 1}^{n}\left|a_i\right|^2\right)
\left(\sum_{i = 1}^{n}\left|b_i\right|^2\right)
\geqslant 
\left|\sum_{i = 1}^{n}a_i\overline{b_i}\right|^2
\]
\end{example}
\end{brown}
\newpage
\section{性质总结}
%均为$n$阶方阵.
\subsection{转置}
\begin{formal}
\begin{proposition}[转置性质总结]\label{prop:转置性质总结}  
$\bm{A}$、
$\bm{B}$均为$n$阶方
阵,$c \in \mathbb{R}$
\begin{enumerate}[label=\textup{(\arabic*)}]
    \item $\left(\bm{A}^{\prime}\right)^{\prime} = \bm{A}$
    \item $\left(\bm{AB}\right)^{\prime} = 
    \bm{B}^{\prime}\bm{A}^{\prime} $
    \item $\left(c\bm{A}\right)^{\prime} = c\bm{A}^{\prime}$
    \item $\left|\bm{A}^{\prime}\right|
    =\left|\bm{A}\right|$
\end{enumerate}
\end{proposition}
\end{formal}
\subsection{逆阵}
\begin{formal}
\begin{proposition}[逆阵性质总结]\label{prop:逆阵性质总结}
$\bm{A}$、
$\bm{B}$均为$n$阶可逆方
阵$\left(
    \left|\bm{A}\right|
    =\left|\bm{B}\right|
    \neq 0
\right)$,$c \in \mathbb{R}$
\begin{enumerate}[label=\textup{(\arabic*)}]
    \item 逆阵存在则唯一
    \item $ \left(\bm{A}^{-1}\right)^{-1}
    = \bm{A} $
    \item $\left(\bm{AB}\right)^{-1}
    = \bm{B}^{-1}\bm{A}^{-1}$
    \item$\left(c\bm{A}\right)^{-1}
    =c^{-1}\bm{A}^{-1}$
    \item $\left|\bm{A}^{-1}\right|
    = \left|\bm{A}\right|^{-1}$
    \item $\left(\bm{A}^{\prime}\right)^{-1}
    = \left(\bm{A}^{-1}\right)^{\prime}$
    \item 整性对可逆阵成立
    \item 消去律对可逆阵成立
\end{enumerate}
\end{proposition}
\end{formal}
\subsection{伴随阵}
\begin{formal}
\begin{proposition}[伴随阵性质总结]\label{prop:伴随阵性质总结}
$\bm{A}$、
$\bm{B}$均为$n$阶方
阵,$c \in \mathbb{R}$,
$n \geqslant 2$
\begin{enumerate}[label=\textup{(\arabic*)}]
    \item 伴随阵必定存在
    \item $\bm{AA}^* = \bm{A}^*\bm{A}
    = \left|\bm{A}\right|\bm{I}_n$
   \item  $\left(\bm{A}^*\right)^*=
    \begin{cases*}
        \left|\bm{A}\right|
        ^{n - 2}\bm{A},n\geqslant 3\\
        \bm{A},n = 2
    \end{cases*}$
    \item $\left(\bm{AB}\right)^* = 
    \bm{B}^*\bm{A}^*$
    \item $\left(c\bm{A}\right)^*=
    c^{n - 1}\bm{A}^*$
    \item $\left|\bm{A}^*\right|
    =\left|\bm{A}\right|^{n - 1}$
    \item $\left(\bm{A}^{\prime}\right)^*
    = \left(\bm{A}^*\right)^{\prime}$
    \item $\left|\bm{A}\right| \neq 0$
        时
    ,$\left(\bm{A}^{-1}\right)^*
    =\left(\bm{A}^*\right)^{-1}$
    \item $\begin{pmatrix}
        \bm{I}_r & \bm{O} \\
        \bm{O} & \bm{O}
    \end{pmatrix}
    ^* = \begin{cases*}
        \bm{O},r \leqslant n - 2
        \\
        \begin{pmatrix}
          \bm{O}  & \bm{O}\\
          0  & 1
        \end{pmatrix},
        r = n -1
        \\
        \bm{I}_n,
        r = n
    \end{cases*}\quad
    \left(0 \leqslant r
    \leqslant n\right)
    $
    \item $r < n$时
$\displaystyle
\left|\begin{pmatrix}
        \bm{I}_r & \bm{O} \\
        \bm{O} & \bm{O}
    \end{pmatrix}
    ^*\right|=0
$
\item $\displaystyle
\left(
    \begin{pmatrix}
        \bm{I}_r & \bm{O} \\
        \bm{O} & \bm{O}
    \end{pmatrix}
    ^*
\right)^*=\begin{cases*}
\begin{pmatrix}
    \bm{I}_r&\bm{O}\\
    \bm{O}&\bm{O}
\end{pmatrix},n=2,r=1\\
\bm{I}_n,r=n\\
\bm{O},\text{其他}
\end{cases*}
$
\end{enumerate}\end{proposition}
\end{formal}
\newpage
\section{一些特别的矩阵}
介绍一些应用广泛的特殊矩阵.
\subsection{标准单位向量}
\begin{formal}
\begin{definition}[标准单位列向量]\label{def:标准单位列向量}
    前文已述,$n$维标准单位列向量指
\[
\bm{e}_1=\begin{pmatrix}
    1\\0\\\vdots\\0
\end{pmatrix},\bm{e}_2=\begin{pmatrix}
    0\\1\\\vdots\\0
\end{pmatrix},\cdots,\bm{e}_n=\begin{pmatrix}
    0\\0\\\vdots\\1
\end{pmatrix}
\]
\end{definition}
\end{formal}
\begin{formal}
\begin{proposition}[标准单位向量的性质]\label{prop:标准单位向量的性质}
另设$\bm{f}_1,\bm{f}_2,\cdots,\bm{f}_m$为$m$维标准单位列向量,则
\begin{enumerate}[label={\textup{(\arabic*)}}]
    \item $\bm{e}_i'\bm{e}_j =0\left(i\neq 0\right),\bm{e}_i'\bm{e}_i=1$
    \item 设$\bm{A}=\left(a_{ij}\right)_{m\times n}$,则$\bm{Ae}_i$是$\bm{A}$的第$i$个列向量,$\bm{f}_i'\bm{A}$是$\bm{A}$的第$i$个行向量
    \item 设$\bm{A}=\left(a_{ij}\right)_{m\times n}$,则$\bm{f}_i'\bm{A}\bm{e}_j=a_{ij}$
    \item 设$\bm{A}=\left(a_{ij}\right)_{m\times n},\bm{B}=\left(b_{ij}\right)_{m\times n}$,则$\bm{A}=\bm{B}$当且仅当$\bm{Ae}_i=\bm{Be}_i\left(\forall 1\leqslant i\leqslant n\right)$,亦当且仅当$\bm{f}_i'\bm{A}=\bm{f}_i'\bm{B}\left(\forall 1\leqslant i\leqslant m\right)$
\end{enumerate}
\end{proposition}
\end{formal}
\subsection{基础矩阵}
\begin{formal}
\begin{definition}[基础矩阵/初级矩阵]\label{def:基础矩阵}
    $n$阶基础矩阵（初级矩阵）是指$n^2$个$n$阶矩阵$\left\{
        \bm{E}_{ij},1\leqslant i,j\leqslant n
    \right\}.$此处$\bm{E}_{ij}$是一个除$\left(i,j\right)$元为$1$外其余元素均为$0$的矩阵即
    \[
    \bm{E}_{ij}=\bm{e}_i\bm{e}_j'
    \]
\end{definition}
\end{formal}
\begin{formal}
\begin{proposition}[基础矩阵的若干性质]
    \,
    \begin{enumerate}[label={\textup{(\arabic*)}}]
        \item 若$j\neq k$,则$\bm{E}_{ij}\bm{E}_{kl}=\bm{O}$
        \item 若$j=k$,则$\bm{E}_{ij}\bm{E}_{kl}=\bm{E}_{il}$
        \item 若$\bm{A}=\left(a_{ij}\right)_{n\times n}$,则
        \[
        \bm{A}=\sum_{i,j=1}^{n}a_{ij}\bm{E}_{ij}
        \]
        \item 若$\bm{A}=\left(a_{ij}\right)_{n\times n}$,则$\bm{E}_{ij}\bm{A}$的第$i$行是$\bm{A}$的第$j$行,
        $\bm{E}_{ij}\bm{A}$的其余行全为零
        \item 若$\bm{A}=\left(a_{ij}\right)_{n
        \times n}$,则$\bm{A}\bm{E}_{ij}$的第$j$列是$\bm{A}$的第$i$列,
        $\bm{A}\bm{E}_{ij}$的其余列全为零
        \item 若$\bm{A}=\left(a_{ij}\right)_{n\times n}$,则$\bm{E}_{ij}\bm{A}\bm{E}_{kl}=a_{jk}\bm{E}_{il}$
    \end{enumerate}
\end{proposition}
\end{formal}
\subsection{循环矩阵}
\begin{formal}
\begin{definition}[循环矩阵定义]\label{def:循环矩阵定义}
    设数域$\mathbb{K}$上的任意$n$个数$a_1,a_2,\cdots,a_n$,则下面的矩阵称为$\mathbb{K}$上的$n$阶循环矩阵
    \[
\bm{A}=\begin{pmatrix}
    a_1&a_2&a_3&\cdots&a_n\\
    a_n&a_1&a_2&\cdots&a_{n-1}\\
    a_{n-1}&a_n&a_1&\cdots&a_{n-2}\\
    \vdots&\vdots&\vdots&&\vdots\\
    a_2&a_3&a_4&\cdots&a_{1}
\end{pmatrix}
    \]
    特别地,$a_2=1,a_1=a_3=\cdots=a_n$的循环矩阵为基础循环矩阵
    \[
    \bm{J}=\begin{pmatrix}
        0&1&0&\cdots&0\\
        0&0&1\cdots&0\\
        \vdots&\vdots&\vdots&&\vdots\\
        0&0&0&\cdots&1\\
        1&0&0&\cdots&0
    \end{pmatrix}
    \]
    任一循环矩阵均可用关于基础矩阵的多项式表出.
\end{definition}
\end{formal}
基础矩阵有非常良好的幂形式,这也是任一循环矩阵可用其的多项式表出的原因,当然,从线性空间的角度也能给出解释.
\begin{green}
\begin{lemma}[基础循环矩阵的幂]\label{lem:基础循环矩阵的幂}
    基础循环矩阵
    \[
    \bm{J}=\begin{pmatrix}
        0&1&0&\cdots&0\\
        0&0&1\cdots&0\\
        \vdots&\vdots&\vdots&&\vdots\\
        0&0&0&\cdots&1\\
        1&0&0&\cdots&0
    \end{pmatrix}
    \]
    的幂次为
    \[
    \bm{J}^k=\begin{pmatrix}
        \bm{O}&\bm{I}_{n-k}\\
        \bm{I}_k&\bm{O}
    \end{pmatrix}
    \left(1\leqslant k\leqslant n\right)
    \]
\end{lemma}
\end{green}
\begin{formal}
\begin{theorem}[循环矩阵关于基础循环矩阵的多项式表示]\label{thm:循环矩阵关于基础循环矩阵的多项式表示}
    利用基础矩阵的幂,容易有
    \[
    \bm{A}=a_1\bm{I}_1+a_2\bm{J}+a_3\bm{J}^2
    +\cdots+a_n\bm{J}^{n-1}
    \]
    令$f\left(x\right)=a_1+a_2x+a_3x^2+\cdots+a_n$是
    $\mathbb{K}$上次数不高于$n-1$的多项式,使得$\bm{A}=f\left(\bm{J}\right)$,
    这就是循环矩阵关于基础循环矩阵的多项式表达.
\end{theorem}
\end{formal}
\begin{formal}
    \begin{theorem}[全体循环矩阵]\label{thm:全体循环矩阵}
        记$C_n\left(\mathbb{K}\right)$为$\mathbb{K}$上的所有
        $n$阶循环矩阵组成的集合,其在矩阵运算
        下是$\mathbb{K}$上的
        $n$维线性空间,一组基为$\left\{
            \bm{I}_n,\bm{J},\cdots,\bm{J}^{n-1}
        \right\}$.
    \end{theorem}
\end{formal}
\begin{formal}
\begin{theorem}[基础循环矩阵的特征值]\label{thm:基础循环矩阵的特征值}
    解特征多项式
    \[
    \left|\lambda \bm{I}_n-\bm{J}\right|=\lambda^n-1
    \]
    即特征值为$n$次单位根
    \[
    \varepsilon_k = \cos\frac{2k\pi}{n}+\mathrm{i}\sin\frac{2k\pi}{n}\left(0\leqslant k\leqslant n - 1\right)
    \]
    其中,特征值$\varepsilon_k$的特征向量为
    \newpage
    \[
    \bm{\alpha}_k=
    \begin{pmatrix}
        1\\\varepsilon_k\\\varepsilon_k^2\\\vdots\\\varepsilon_k^{n-1}
    \end{pmatrix}
    \]
\end{theorem}
\end{formal}
\subsection{幂零Jordan块}
\begin{formal}
\begin{definition}[幂零Jordan块]\label{def:幂零Jordan块}
    跟基础循环矩阵很像
\[
\bm{A}=\begin{pmatrix}
        0&1&0&\cdots&0\\
        0&0&1\cdots&0\\
        \vdots&\vdots&\vdots&&\vdots\\
        0&0&0&\cdots&1\\
        0&0&0&\cdots&0
    \end{pmatrix}
\]
其幂次
\[
\bm{A}^k=\begin{pmatrix}
    \bm{O}&\bm{I}_{n-k}\\
    \bm{O}&\bm{O}
\end{pmatrix}\left(1\leqslant k\leqslant 
n\right)
\]
\end{definition}
\end{formal}
\subsection{多项式的友阵与Frobenius块}
\begin{formal}
\begin{definition}[定义]\label{def:多项式的友阵与Frobenius块}
    设首一多项式$f\left(x\right)=x^n+a_1x^{n-1}+\cdots+a_{n-1}x+a_n$,其友阵为
    \[
    \bm{C}\left(f\left(x\right)\right)=\begin{pmatrix}
        0&0&\cdots&0&-a_n\\
        1&0&\cdots&0&-a_{n-1}\\
        0&1&\cdots&0&-a_{n-2}\\
        \vdots&\vdots&&\vdots&\vdots\\
        0&0&\cdots&1&-a_1
    \end{pmatrix}
    \]
    其转置$\bm{F}\left(f\left(x\right)\right)=
    \bm{C}\left(f\left(x\right)\right)$称为$\textup{Frobenius}$块.
\end{definition}
\end{formal}
\newpage
\begin{formal}
\begin{proposition}[友阵的性质]\label{prop:友阵的性质}
    设$\bm{e}_i$是标准单位列向量,则\[
    \bm{C}\left(f\left(x\right)\right)\bm{e}_i=e_{i+1}\left(\forall 1\leqslant i\leqslant n-1\right),
    \bm{C}\left(f\left(x\right)\right)\bm{e}_n=-\sum_{i=1}^{n}a_{n-i+1}\bm{e}_i
    \]
\end{proposition}
\end{formal}
\begin{formal}
\begin{theorem}[友阵与原多项式关系]\label{thm:友阵与原多项式关系}
    \[
    \left|x\bm{I}_n-\bm{C}\left(f\left(x\right)\right)\right|=f\left(x\right)
    \]
\end{theorem}
\begin{Proof}
考虑行列式
\[
F_n=\begin{vmatrix}
    \lambda&0&0&\cdots&0&a_n\\
    -1&\lambda&0&\cdots&0&a_{n-1}\\
    0&-1&\lambda&\cdots&0&a_{n-2}\\
    \vdots&\vdots&\vdots&&\vdots&\vdots\\
    0&0&0&\cdots&\lambda&a_2\\
    0&0&0&\cdots&-1&\lambda+a_1
\end{vmatrix}
\]
按第一行展开,注意到：$a_n$的余子式是上三角,$\lambda$的余子式是更低一阶的类似行列式,即\[
F_n=\lambda F_{n-1}+\left(-1\right)^{n+1}\left(-1\right)^{n-1}a_n=\lambda F_{n-1}+a_n
\]
因此\[
F_n=\lambda^n+a_1\lambda^{n-1}+a_2\lambda^{n-2}+\cdots+a_n\qedhere
\]
\end{Proof}
\end{formal}
\begin{red}
\begin{remark}
    \textup{Frobenius}块与友阵将在后面的\cref{lemma:友阵}有理标准型学习中出现.
\end{remark}
\end{red}
\subsection{对称阵与反对称阵}
\begin{formal}
\begin{proposition}[对对称阵与反对称阵的刻画]
    \,
    \begin{enumerate}[label={\textup{(\arabic*)}}]
        \item 设$n$阶对称阵$\bm{A}$,则$\bm{A}$是零矩阵等价于$\forall \bm{\alpha}\in\mathbb{R}^n$
        \[
        \bm{\alpha}'\bm{A\alpha}=0
        \]
        \item 设$n$阶方阵$\bm{A}$,则$\bm{A}$是反对称阵等价于$\forall \bm{\alpha}\in\mathbb{R}^n$
        \[
        \bm{\alpha}'\bm{A\alpha}=0
        \]
    \end{enumerate}
\end{proposition}
\end{formal}